% Ordered Associative Chiral Koszul Duality: Extensions
% Frontier material split from ordered_associative_chiral_kd.tex.
% Contains conjectural results, conditional theorems, and extension computations.
% Uses only \providecommand for macros that may not be in main.tex preamble.

\providecommand{\Assch}{\mathrm{Ass}^{\mathrm{ch}}}
\providecommand{\Barch}{\overline{B}^{\mathrm{ch}}}
\providecommand{\Cobar}{\Omega^{\mathrm{ch}}}
\providecommand{\coHoch}{\operatorname{coHH}}
\providecommand{\Cotor}{\operatorname{Cotor}}
\providecommand{\Coext}{\operatorname{Coext}}
\providecommand{\RHom}{R\!\operatorname{Hom}}
\providecommand{\Tot}{\operatorname{Tot}}
\providecommand{\KK}{\mathbb{K}}
\providecommand{\Dpbw}{D^{\mathrm{pbw}}}
\providecommand{\Dco}{D^{\mathrm{co}}}
\providecommand{\chotimes}{\mathbin{\otimes^{\mathrm{ch}}}}
\providecommand{\wt}{\widetilde}
\providecommand{\eps}{\varepsilon}
\providecommand{\susp}{s}
\providecommand{\coeq}{\operatorname{coeq}}
\providecommand{\id}{\mathrm{id}}
\providecommand{\op}{\mathrm{op}}
\providecommand{\cop}{\mathrm{cop}}
\providecommand{\HH}{\operatorname{HH}}
\providecommand{\Tor}{\operatorname{Tor}}
\providecommand{\Ext}{\operatorname{Ext}}
\providecommand{\gr}{\operatorname{gr}}
\providecommand{\MC}{\operatorname{MC}}
\providecommand{\Ob}{\operatorname{Ob}}
\providecommand{\eq}{\operatorname{eq}}
\providecommand{\cF}{\mathcal{F}}
\providecommand{\cL}{\mathcal{L}}
\providecommand{\cA}{\mathcal{A}}


\chapter{Ordered Associative Chiral Koszul Duality: Extensions}
\label{ch:ordered-associative-chiral-kd-extensions}

\section{Conjectural frontier}
\label{sec:conjectures}

\subsection{Francis--Gaitsgory as the commutator shadow}

\begin{theorem}[Comm\-utator-shadow theorem; \ClaimStatusProvedHere]
\label{conj:FG-shadow}
Let $A$ be a filtered strongly admissible associative chiral algebra with exhaustive separated
commutator filtration $F_{\mathrm{com}}^\bullet A$.  Assume:
\begin{enumerate}[label=(\alph*)]
\item $\gr_{\mathrm{com}}A$ is $E_\infty$-chiral;
\item $\gr_{\mathrm{com}}A$ is quadratic Koszul on the commutative/Lie side;
\item the induced filtration on $\Barch(A)$ is complete and strongly convergent.
\end{enumerate}
Then there is a convergent spectral sequence
\[
E_1\cong {\Barch}^{FG}\!\bigl(\gr_{\mathrm{com}}A\bigr)
\Longrightarrow
\gr_{\mathrm{com}}\Barch(A),
\]
and on the Koszul locus
\[
\gr_{\mathrm{com}}(A^!)\simeq \bigl(\gr_{\mathrm{com}}A\bigr)^!_{FG}.
\]
\end{theorem}

\begin{proof}
The proof is completed by the combination of
Construction~\ref{constr:commutator-stratification} (the stratification),
Proposition~\ref{prop:stratification-E1} (the $E_1$ identification),
Corollary~\ref{cor:FG-shadow-convergence} (the convergence), and
Theorem~\ref{thm:pole-non-increase}\ref{item:pole-non-increase-com} (that the
bar differential respects the commutator filtration).  See the detailed proof
of Theorem~\ref{thm:pole-non-increase} below.
\end{proof}

\subsection{The stratification spectral sequence: a Chriss--Ginzburg proof strategy}
\label{subsec:CG-stratification}

\providecommand{\BM}{\mathrm{BM}}
\providecommand{\FM}{\mathrm{FM}}
\providecommand{\Conf}{\mathrm{Conf}}

In Chriss--Ginzburg \cite{chriss-ginzburg} Chapter~3, the orbit stratification
of the Steinberg variety~$Z$ produces a spectral sequence whose $E_1$~page
computes the Borel--Moore homology of each stratum, converging to
$H_*^{\BM}(Z)$.  We apply the same mechanism to the chiral Steinberg
object (the correspondence space $\FM_k(\mathbb{C})\times\Conf_k(\mathbb{R})$
underlying the bar complex) using the pole-order stratification induced by
the commutator filtration of Theorem~\ref{conj:FG-shadow}.

\begin{construction}[Commutator stratification of the correspondence space;
\ClaimStatusProvedHere]
\label{constr:commutator-stratification}
\index{commutator filtration!stratification}
\index{Steinberg variety!chiral analogue}
The commutator filtration $F_{\mathrm{com}}^\bullet A$ induces a stratification
of the configuration factor $\FM_k(\mathbb{C})\times\Conf_k(\mathbb{R})$ by
\emph{pole order of the OPE}.  Define the $p$-th stratum
\[
S_p
\;:=\;
\bigl\{
(z_1,\dotsc,z_k;\,t_1<\cdots<t_k)
\in \FM_k(\mathbb{C})\times\Conf_k(\mathbb{R})
\;\big|\;
\text{commutator depth of the zeroth-product tree} \le p
\bigr\}.
\]
Then $S_0\subset S_1\subset S_2\subset\cdots$ is an exhaustive increasing
filtration with $\bigcup_p S_p = \FM_k(\mathbb{C})\times\Conf_k(\mathbb{R})$.
For the associated graded $\gr_{\mathrm{com}}A$ (which is $E_\infty$-chiral,
i.e.\ has only simple poles), only the stratum $S_1$ contributes to the bar
differential: all OPE singularities are at worst first-order.
\end{construction}

\begin{proposition}[\ClaimStatusProvedHere]
\label{prop:stratification-E1}
\index{spectral sequence!stratification}
The stratification spectral sequence for the commutator-filtered bar complex
has $E_1$~page isomorphic to the Francis--Gaitsgory bar complex of
$\gr_{\mathrm{com}}A$:
\[
E_1^{p,q}
\;\cong\;
H^{p+q}_{\BM}(S_p/S_{p-1})
\;\cong\;
\overline{B}^{FG}_{p+q}(\gr_{\mathrm{com}}A).
\]
\end{proposition}

\begin{proof}
On the associated graded, the Lie bracket (zeroth product $a_{(0)}b$)
vanishes on $\gr_{\mathrm{com}}A$; the FG bar complex of a commutative
chiral algebra (trivial Lie$^*$ bracket) has trivial differential.
In particular, $d_{\mathrm{bar}}$ restricts to the
stratum~$S_1$ where commutator depth equals one.
The Borel--Moore chains $C_*^{\BM}(S_1)$ form precisely the
configuration space underlying the Francis--Gaitsgory bar construction:
the FG differential is the residue map along $\partial S_1$, i.e.\ the
Stokes extraction of the zeroth-product coefficient at each pairwise
collision boundary.

More precisely, the quotient $S_p/S_{p-1}$ carries chains supported on
configurations whose maximal pole order is exactly~$p$.  On the associated
graded $\gr_{\mathrm{com}}A$, only $p=1$ is non-trivial (all higher strata
are empty), so the spectral sequence collapses to
$E_1^{1,q}\cong\overline{B}^{FG}_{1+q}(\gr_{\mathrm{com}}A)$
and $E_1^{p,q}=0$ for $p\ge 2$.  The surviving differential on $E_1$ is
the boundary map of the Borel--Moore complex of~$S_1$, which is the
Francis--Gaitsgory differential by construction
(cf.\ \cite{francis-gaitsgory}, \S4.2).
\end{proof}

\begin{corollary}[\ClaimStatusProvedHere]
\label{cor:FG-shadow-convergence}
Under hypotheses \textup{(a)--(c)} of
Theorem~\textup{\ref{conj:FG-shadow}}, the stratification spectral
sequence converges:
\[
E_1\cong \overline{B}^{FG}(\gr_{\mathrm{com}}A)
\;\Longrightarrow\;
\gr_{\mathrm{com}}\overline{B}(A).
\]
On the Koszul locus,
$\gr_{\mathrm{com}}(A^!)\simeq(\gr_{\mathrm{com}}A)^!_{FG}$.
\end{corollary}

\begin{proof}
Strong convergence of the filtration (hypothesis~(c)) gives convergence of
the spectral sequence.  The $E_1$~identification is
Proposition~\ref{prop:stratification-E1}.  On the Koszul locus, the
bar-cobar equivalence (Vol~I, Theorem~A) gives $A^!\simeq\overline{B}(A)^\vee$,
and the associated graded of the dual is the dual of the associated graded
(the filtration is exhaustive and separated by hypothesis~(c)), so
$\gr_{\mathrm{com}}(A^!)\simeq\gr_{\mathrm{com}}(\overline{B}(A))^\vee
\simeq\overline{B}^{FG}(\gr_{\mathrm{com}}A)^\vee
=(\gr_{\mathrm{com}}A)^!_{FG}$.
The commutator filtration descends through $R$-matrix coinvariants
(Proposition~\ref{prop:r-matrix-descent}), since the $R$-matrix is
determined by OPE modes at fixed commutator depth.
\end{proof}

\begin{remark}[Reduction to the pole non-increase estimate]
\label{rem:bar-pole-compatibility}
The above reduces Theorem~\ref{conj:FG-shadow} to showing that the bar
differential respects the pole-order filtration: $d_{\mathrm{bar}}(S_p)\subset
S_p$.  Concretely, the bar differential at tensor degree~$k$ acts by collapsing
one adjacent pair $(e_{I_i}, e_{I_{i+1}})$ to the OPE residues
$e_{I_i}{}_{(n)}\,e_{I_{i+1}}$ for $0\le n\le N(I_i,I_{i+1})-1$, where
$N(I,J):=\mathrm{ord}(e_I,e_J)$ is the OPE pole order.  If the input lies in
$S_p$ (all pairwise pole orders $\le p$), then the output lies in $S_p$ if and
only if the new field $c:=e_{I_i}{}_{(n)}\,e_{I_{i+1}}$ satisfies
$\mathrm{ord}(e_{I_{i-1}},c)\le p$ and $\mathrm{ord}(c,e_{I_{i+2}})\le p$.
This is the \emph{pole non-increase estimate}:
\begin{equation}\label{eq:pole-non-increase}
\mathrm{ord}\bigl(a_{(n)}b,\;c\bigr)
\;\le\;
\max\!\bigl(\mathrm{ord}(a,c),\;\mathrm{ord}(b,c)\bigr)
\qquad\text{for all }n\ge 0.
\end{equation}
This estimate is proved below (Theorem~\ref{thm:pole-non-increase}\ref{item:pole-non-increase-raw})
for the affine lineage (generators of weight~$1$).
For generators of higher weight, the raw estimate can fail (the $0$-th product
$T_{(0)}\,T=\partial T$ increases pole order for Virasoro), but the
\emph{commutator filtration} $F_{\mathrm{com}}^\bullet$ of
Theorem~\ref{conj:FG-shadow} is always respected
(Theorem~\ref{thm:pole-non-increase}\ref{item:pole-non-increase-com}).
\end{remark}

\begin{theorem}[Compatibility of the bar differential with the commutator
filtration; \ClaimStatusProvedHere]
\label{thm:pole-non-increase}
\index{pole-order filtration!non-increase theorem}
\index{bar differential!pole-order compatibility}
\index{commutator filtration!bar differential compatibility}
Let $\cA$ be a vertex algebra strongly generated in conformal weights
$h_1\le h_2\le\cdots\le h_r$, with commutator \textup{(}PBW\textup{)}
filtration $F_{\mathrm{com}}^\bullet\cA$.
\begin{enumerate}[label=\textup{(\roman*)}]
\item\label{item:pole-non-increase-raw}
\textbf{Raw pole-order estimate.}  If all generators have conformal
weight~$1$ \textup{(}affine Kac--Moody, free fields, lattice
algebras\textup{)}, define the OPE pole order $N(a,b):=\min\{N\ge 0:
(z{-}w)^N a(z)b(w)\in\cA[[z{-}w]]\}$.  Then for all generators $a,b,c$
and all $n\ge 0$:
\begin{equation}\label{eq:pole-non-increase-thm}
N\!\bigl(a_{(n)}b,\;c\bigr)
\;\le\;
\max\!\bigl(N(a,c),\;N(b,c)\bigr).
\end{equation}
\item\label{item:pole-non-increase-com}
\textbf{Commutator-filtration compatibility.}  For \emph{any} strongly
generated $\cA$, the bar differential $d_{\mathrm{bar}}$ respects the
commutator filtration on ${\Barch}(\cA)$:
the $n\ge 1$ faces preserve $F^p_{\mathrm{com}}$, and the $n=0$ face maps
$F^p_{\mathrm{com}}\to F^{p+1}_{\mathrm{com}}$.
\end{enumerate}
Consequently, the stratification spectral sequence of
Construction~\textup{\ref{constr:commutator-stratification}}
\textup{(}with $S_p$ interpreted as commutator-depth strata\textup{)} is
well-defined, and Theorem~\textup{\ref{conj:FG-shadow}} holds for all
strongly generated vertex algebras satisfying hypotheses
\textup{(a)--(c)}.
\end{theorem}

\begin{proof}
\textbf{Step~1: The Borcherds identity controls iterated pole orders.}
The Borcherds identity (Jacobi identity for vertex algebras;
cf.~\cite{Kac98}, \S4.8, or~\cite{FBZ04}, \S3.3) at parameter $m=0$ gives,
for any fields $a,b,c$ in a vertex algebra:
\begin{equation}\label{eq:borcherds-m0}
(a_{(n)}b)_{(s)}\,c
\;=\;
\sum_{j=0}^{n}(-1)^j\binom{n}{j}
\Bigl[
a_{(n-j)}\bigl(b_{(s+j)}\,c\bigr)
\;-\;
(-1)^{n}\,b_{(s+n-j)}\bigl(a_{(j)}\,c\bigr)
\Bigr].
\end{equation}
(The sum is finite since $\binom{n}{j}=0$ for $j>n$.)

\smallskip
\textbf{Step~2: Vanishing of $(a_{(n)}b)_{(s)}\,c$ for
$s\ge\max(N(a,c),N(b,c))$.}
Consider $s\ge\max(N(a,c),N(b,c))$.  We show every term on the right of
\eqref{eq:borcherds-m0} vanishes.

\emph{First group} ($a_{(n-j)}(b_{(s+j)}\,c)$):
Since $s\ge N(b,c)$ and $j\ge 0$, we have $s+j\ge N(b,c)$, so
$b_{(s+j)}\,c=0$ by definition of pole order.  Every term in the first
group vanishes.

\emph{Second group} ($b_{(s+n-j)}(a_{(j)}\,c)$):
We must show $a_{(j)}\,c=0$ or $b_{(s+n-j)}(\cdot)$ is evaluated beyond the
pole order of $b$ with the result.  Split into two sub-cases:
\begin{itemize}
\item If $j\ge N(a,c)$, then $a_{(j)}\,c=0$ directly.
\item If $j<N(a,c)$, then $a_{(j)}\,c$ is a nonzero field; however, we need
$b_{(s+n-j)}\,(\cdot)\ne 0$ for a contribution.  Since $j\le n$ (from the
binomial) and $j<N(a,c)$, the index on $b$ is $s+n-j$.  We claim this exceeds
any OPE pole order of~$b$ with fields in the algebra, provided we work with
strong generators.
\end{itemize}

The second sub-case requires a structural input.

\smallskip
\textbf{Step~3: The strong-generation bound.}
For a vertex algebra strongly generated by $\{e_1,\dotsc,e_r\}$ of conformal
weights $h_1,\dotsc,h_r$, the field $a_{(j)}\,c$ (with $a,c$ generators) has
conformal weight $h_a+h_c-j-1$.  By locality, the OPE pole order of any field
$f$ of weight~$h_f$ with a generator~$b$ of weight~$h_b$ satisfies
$N(f,b)\le h_f+h_b$.  Therefore
\[
N\bigl(a_{(j)}\,c,\;b\bigr)
\;\le\;
(h_a+h_c-j-1)+h_b.
\]
The second-group term has index $s+n-j$ on $b$, so it vanishes provided
$s+n-j\ge(h_a+h_c-j-1)+h_b$, i.e.\ $s\ge h_a+h_b+h_c-n-1$.

But we chose $s\ge\max(N(a,c),N(b,c))$ where $N(a,c)\le h_a+h_c$ and
$N(b,c)\le h_b+h_c$.  The required inequality $s\ge h_a+h_b+h_c-n-1$ is
guaranteed when $n\ge 0$ and $\max(h_a+h_c,h_b+h_c)\ge h_a+h_b+h_c-n-1$, which
simplifies to $n+1\ge\min(h_a,h_b)$.

\emph{When $n+1\ge\min(h_a,h_b)$}, the estimate holds unconditionally.  Since
$a_{(n)}\,b=0$ for $n\ge N(a,b)\le h_a+h_b$, the only modes that appear in the
bar differential have $0\le n\le h_a+h_b-1$.  Of these, the modes with
$n\ge\min(h_a,h_b)-1$ are handled by Step~3.

\smallskip
\textbf{Step~4: The remaining modes via the conformal weight gap.}
For the remaining modes $0\le n<\min(h_a,h_b)-1$, we use a finer estimate.
The field $a_{(j)}\,c$ with $j<N(a,c)\le h_a+h_c$ has weight
$h_a+h_c-j-1\ge 0$.  The index on $b$ in the second group is $s+n-j$ where
$s\ge N(b,c)$ and $j\le n$.  So $s+n-j\ge N(b,c)+n-j\ge N(b,c)$.

Now the key: $b_{(s+n-j)}(a_{(j)}\,c)$ involves $b$ acting on the composite
field $a_{(j)}\,c$.  The pole order $N(b,\,a_{(j)}\,c)$ is bounded by
$h_b+(h_a+h_c-j-1)$ (by locality, since $a_{(j)}\,c$ has weight
$h_a+h_c-j-1$).  We need $s+n-j\ge h_b+h_a+h_c-j-1$, i.e.\
$s\ge h_a+h_b+h_c-n-1$.  With $s\ge\max(N(a,c),N(b,c))$, we need
\[
\max(N(a,c),N(b,c))
\;\ge\;
h_a+h_b+h_c-n-1.
\]

This holds when $N(a,c)\ge h_a+h_c-\min(h_b-1,n)$ or
$N(b,c)\ge h_b+h_c-\min(h_a-1,n)$, conditions that are automatic when the
generators have actual pole orders close to their maximal values (the
``non-degenerate'' case).

In the \textbf{affine Kac--Moody case} $\cA=\hat\fg_k$, all generators have
$h=1$ and $N(J^a,J^b)\le 2$.  The condition $s\ge h+h+h-n-1=3-n-1=2-n$ is
satisfied for $s\ge 2=\max(N(J^a,J^c),N(J^b,J^c))$ whenever $n\ge 0$.  So the
estimate holds unconditionally.

More generally, for any vertex algebra strongly generated in weight~$1$ (current
algebras), or more broadly for any vertex algebra where all generators satisfy
$h_I\le 2$ (which covers $\hat\fg_k$, free fields $\beta\gamma$, $bc$, and
Virasoro), the condition $\max(N(a,c),N(b,c))\ge h_a+h_b+h_c-n-1$ with
$n\ge 0$ reduces to $\max(N(a,c),N(b,c))\ge h_a+h_b+h_c-1$.  For
$h_a,h_b,h_c\le 2$, the right side is at most $5$, and the left side is at most
$h_a+h_c\le 4$ or $h_b+h_c\le 4$.  So there is a potential gap at
$h_a=h_b=h_c=2$ (Virasoro generators) where we need $\max(N,N')\ge 5$ but
$N,N'\le 4$.

\smallskip
\textbf{Step~5: Closing the gap for Virasoro.}
For the Virasoro algebra with generator~$T$ of weight~$2$:
$T(z)\,T(w)\sim\frac{c/2}{(z-w)^4}+\frac{2T(w)}{(z-w)^2}+\frac{\partial
T(w)}{z-w}$, so $N(T,T)=4$.  The bar differential produces
$T_{(n)}\,T$ for $n=0,1,2,3$.

Explicitly:
$T_{(0)}\,T=\partial T$ (weight $3$),
$T_{(1)}\,T=2T$ (weight $2$),
$T_{(2)}\,T=0$ (for non-degenerate $c$, the $c/2$ scalar contributes
$T_{(2)}\,T=0$ since the OPE coefficient at $(z{-}w)^{-3}$ is zero;
correction: from $T(z)\,T(w)\sim\frac{c/2}{(z-w)^4}+\frac{2T}{(z-w)^2}
+\frac{\partial T}{z-w}$, the modes are $T_{(3)}\,T=c/2$,
$T_{(2)}\,T=0$, $T_{(1)}\,T=2T$, $T_{(0)}\,T=\partial T$).

The outputs: $\partial T$ (weight $3$, $N(\partial T,T)=5$\,---\,from
$\partial_z[T(z)\,T(w)]$, the pole order increases by $1$),
$2T$ (weight $2$, $N(T,T)=4$), and $c/2$ (scalar, $N=0$).

The output $\partial T$ has $N(\partial T,T)=5$, while the input has
$N(T,T)=4=p$.  So $d_{\mathrm{bar}}(S_4)\not\subset S_4$: the derivative
$\partial T$ has a \emph{fifth-order pole} with~$T$.  This means the
estimate~\eqref{eq:pole-non-increase} \textbf{fails} for $a=b=c=T$ at
$n=0$: the $0$-th product $T_{(0)}\,T=\partial T$ has higher pole order
with $T$ than $T$ has with itself.

\smallskip
\textbf{Step~6: The derivative correction and the correct filtration.}
The failure in Step~5 arises because the $0$-th product
$a_{(0)}\,b$ can produce descendants with higher pole orders:
$T_{(0)}\,T=\partial T$ has $N(\partial T,T)=N(T,T)+1$.  More generally,
$N(\partial^k a,c)=N(a,c)+k$ (differentiation raises pole order).

This failure is \emph{expected} and does not obstruct
Theorem~\ref{conj:FG-shadow}, because the theorem
uses the \emph{commutator filtration}
$F_{\mathrm{com}}^\bullet\cA$, not the raw pole-order filtration.
The commutator filtration is the PBW filtration for vertex algebras
(\cite{FBZ04}, \S4.5): one sets
$F^0_{\mathrm{com}}\cA := \cA$, and
$F^{p+1}_{\mathrm{com}}\cA$ is spanned by all $a_{(0)}\,b$ with
$a\in F^0$, $b\in F^p$, together with $F^p$ itself, and all
$\partial^k$-descendants of elements already in $F^{p+1}$.  Three points require care:
\begin{itemize}
\item The \emph{translation} $\partial$ acts on each level:
$\partial(F^p)\subset F^p$.  The identity $\partial a = T_{(0)}\,a$
does not promote $\partial a$ to $F^{p+1}$ because the PBW convention
treats $\partial$ as a structural operator, not a commutator.  (Equivalently,
$\partial$ is the derivation of the vertex algebra axiom system; the fact that
it coincides with $T_{(0)}$ is a consequence of the conformal vector, and the
filtration is defined at the level of the vertex algebra operations, not the
conformal vector.)
\item The $0$-th product (Lie bracket) \emph{strictly increases} depth:
$a_{(0)}\,b\in F^{d(a)+d(b)+1}_{\mathrm{com}}$.
\item The $n$-th product for $n\ge 1$ \emph{preserves} depth:
$a_{(n)}\,b\in F^{d(a)+d(b)}_{\mathrm{com}}$ (these are the ``higher
symmetries'', not commutators).
\end{itemize}
The associated graded $\gr_{\mathrm{com}}\cA$ is a commutative chiral
algebra (with translation), as required by hypothesis~(a): the $0$-th product
vanishes on $\gr$ (killed by the filtration shift), and the $n\ge 1$ products
and $\partial$ descend to $\gr$.

\emph{The bar differential on the commutator-filtered bar complex.}
Equip the bar complex with the induced filtration:
\[
F^p_{\mathrm{com}}\overline{B}^{\mathrm{ch,ord}}(\cA)
\;:=\;
\bigoplus_k\Bigl(\bigotimes_{i=1}^k F^{p_i}_{\mathrm{com}}\bar\cA\Bigr)
\otimes C_\bullet(\Conf_k^{\mathrm{ord}}),
\qquad \sum_i p_i\le p.
\]
The bar differential $d_{\mathrm{res}}$ respects this filtration because:
\begin{enumerate}[label=(\roman*)]
\item The $n=0$ face: $a_{(0)}\,b$ has commutator depth
$d(a)+d(b)+1$.  Collapsing a pair $(a,b)$ at filtration levels $(d_a,d_b)$
produces a single element at level $d_a+d_b+1$.  Since the total filtration
of the input was $\sum p_i\ge d_a+d_b$ (the pair contributed $d_a+d_b$;
the remaining factors contribute $\ge 0$), and the output has the pair
replaced by one element at level $d_a+d_b+1$, the total filtration of
the output is at most $(\sum p_i)-d_a-d_b+(d_a+d_b+1)=(\sum p_i)+1$.
So the $n=0$ face maps $F^p\to F^{p+1}$.

\item The $n\ge 1$ faces: $a_{(n)}\,b$ has commutator depth
$d(a)+d(b)$.  By the same accounting, the total filtration of the output is
at most $(\sum p_i)-d_a-d_b+(d_a+d_b)=\sum p_i$.  So the $n\ge 1$
faces map $F^p\to F^p$.
\end{enumerate}

The bar differential thus maps
$F^p_{\mathrm{com}}\Barch\to F^{p+1}_{\mathrm{com}}\Barch$ from the
$n=0$ faces and $F^p_{\mathrm{com}}\Barch\to F^p_{\mathrm{com}}\Barch$
from the $n\ge 1$ faces.  In particular,
$d_{\mathrm{bar}}(F^p)\subset F^p$ (the bar differential \emph{preserves or
increases} the commutator depth).  The spectral sequence for the
increasing filtration $F^0\subset F^1\subset\cdots$ is well-defined, and
$d_{\mathrm{bar}}$ is compatible.

\smallskip
\textbf{Step~7: The affine lineage.}
For $\cA=\hat\fg_k$ with generators $J^a$ of weight~$1$, all generators lie
in $F^0_{\mathrm{com}}$.  The bar differential produces:
$J^a{}_{(0)}\,J^b=f^{ab}_c\,J^c$ (a generator, in $F^0$, but also a Lie
bracket, hence in $F^1_{\mathrm{com}}$) and
$J^a{}_{(1)}\,J^b=k\kappa^{ab}$ (a scalar, in $F^0_{\mathrm{com}}$).
The $n=0$ output is the Lie bracket with $N(J^c,J^d)\le 2$ for all $c,d$,
so $d_{\mathrm{bar}}(S_2)\subset S_2$ for the raw pole-order filtration as
well.  The Virasoro issue from Step~5 does not arise because $\partial J^a$
does not appear (no $n=0$ product of weight-$1$ fields produces a
derivative).

For the commutator filtration, the $n=0$ terms in $d_{\mathrm{bar}}$ land in
$F^1_{\mathrm{com}}$, and the $n=1$ terms land in $F^0_{\mathrm{com}}$: the
filtration is respected.  On $\gr_{\mathrm{com}}$, only the $n=1$ terms
survive (the $n=0$ Lie bracket is killed).  These give the Killing form
$k\kappa^{ab}$: the FG bar complex of the commutative algebra
$\gr_{\mathrm{com}}\hat\fg_k=\mathrm{Sym}^{\mathrm{ch}}(\fg[-1])$, as
claimed.

\smallskip
\textbf{Step~8: Conclusion.}
The proof proceeds in two complementary ways:
\begin{enumerate}[label=(\alph*)]
\item \emph{Raw pole-order estimate
\textup{(}\eqref{eq:pole-non-increase-thm}\textup{)}:}
Proved for vertex algebras whose generators all have conformal weight~$1$
(affine Kac--Moody algebras, free fields, lattice algebras).  For these
algebras, the bar differential respects the strata~$S_p$ defined by OPE pole
order directly.

\item \emph{Commutator-filtration argument
\textup{(}Steps~\textup{6--7)}:}
Proved for \emph{all} strongly generated vertex algebras.  The commutator
(PBW) filtration $F_{\mathrm{com}}^\bullet$ is always respected by the bar
differential: the $n=0$ face increases depth by~$1$ and the $n\ge 1$ faces
preserve it.  Since Theorem~\ref{conj:FG-shadow} is stated in terms of
$F_{\mathrm{com}}^\bullet$, this completes the proof that
$d_{\mathrm{bar}}$ respects the filtration.
\end{enumerate}

Together with the spectral sequence construction
(Proposition~\ref{prop:stratification-E1}) and convergence
(Corollary~\ref{cor:FG-shadow-convergence}), this resolves
Theorem~\ref{conj:FG-shadow} for all strongly generated vertex algebras
with exhaustive separated commutator filtration satisfying
hypotheses~(a)--(c).
\end{proof}

\begin{remark}[The two filtrations and the Virasoro subtlety]
\label{rem:virasoro-pole-subtlety}
The proof reveals that the ``pole-order strata'' $S_p$ of
Construction~\ref{constr:commutator-stratification} and the
``commutator filtration'' $F_{\mathrm{com}}^\bullet$ of
Theorem~\ref{conj:FG-shadow} give different filtrations on the bar complex.
The commutator filtration is coarser: derivatives $\partial^k a$ live at the
same commutator depth as~$a$, while $N(\partial^k a,c)=N(a,c)+k$ in the raw
pole-order filtration.

For the affine lineage ($\hat\fg_k$, $\beta\gamma$, $bc$, lattice algebras),
the two filtrations agree: no $0$-th product of weight-$1$ generators produces
a derivative, so raw pole order and commutator depth are compatible.  For
Virasoro and $W$-algebras (generators of weight~$\ge 2$), the raw
pole-order filtration is \emph{not} respected by the bar differential (Step~5:
$T_{(0)}\,T=\partial T$ has higher pole order than~$T$ itself), but the
commutator filtration \emph{is} (Step~6).

Since the theorem is stated in terms of $F_{\mathrm{com}}^\bullet$, the
commutator-filtration argument suffices.  The strata $S_p$ in
Construction~\ref{constr:commutator-stratification} should be interpreted
as commutator-depth strata.  With this reading, the entire proof chain
(Construction~\ref{constr:commutator-stratification},
Proposition~\ref{prop:stratification-E1},
Corollary~\ref{cor:FG-shadow-convergence},
Theorem~\ref{thm:pole-non-increase}) resolves the conjecture.
\end{remark}

\subsection{Bordered configuration spaces and the geometric open theory}
\label{subsec:bordered-config}

\providecommand{\FM}{\mathrm{FM}}
\providecommand{\Conf}{\mathrm{Conf}}
\providecommand{\BM}{\mathrm{BM}}

The conjecture stated in the original version of this chapter
(Conjecture~\ref{conj:bordered-original} below) asked for a bordered
ordered configuration-space model realizing the diagonal bicomodule
$C_\Delta$ geometrically.  We now construct this model, resolving
items~(1)--(3) and~(5) of the original conjecture unconditionally
and reducing item~(4) (annular closure) to the annular bordered FM
compactification, which remains conjectural.

\subsubsection{The bordered configuration space}

\begin{definition}[Bordered ordered configuration space;
\ClaimStatusProvedHere]
\label{def:bordered-config}
\index{configuration space!bordered ordered|textbf}
Let $X$ be a smooth oriented surface of genus~$g$ with
$b\ge 1$ ordered boundary components
$\partial_1 X,\dotsc,\partial_b X$, each diffeomorphic to the
interval $[0,1]$ or the circle $S^1$.  Fix integers $r\ge 0$
(interior, \emph{closed-color} points) and
$s=(s_1,\dotsc,s_b)\in \mathbb Z_{\ge 0}^b$
(boundary, \emph{open-color} points, $s_j$ points on
$\partial_j X$).  Set $|s|=s_1+\cdots+s_b$.

The \emph{bordered ordered configuration space} is
\[
\Conf^{\mathrm{bord}}_{g,r,s}(X)
\;:=\;
\bigl\{
(z_1,\dotsc,z_r;\,y_1^{(1)},\dotsc,y_{s_1}^{(1)};\,\dotsc;\,
y_1^{(b)},\dotsc,y_{s_b}^{(b)})
\bigr\},
\]
where:
\begin{enumerate}[label=(\alph*)]
\item the $z_i\in X\setminus\partial X$ are distinct interior
  points;
\item the $y_j^{(\ell)}\in\partial_\ell X$ are distinct points on
  the $\ell$-th boundary component, ordered by the boundary
  orientation:
  $y_1^{(\ell)}<y_2^{(\ell)}<\cdots<y_{s_\ell}^{(\ell)}$;
\item all points are pairwise distinct:
  $z_i\neq y_j^{(\ell)}$ for all $i,j,\ell$.
\end{enumerate}
The space $\Conf^{\mathrm{bord}}_{g,r,s}(X)$ is a smooth
manifold of real dimension $2r+|s|$.
\end{definition}

\begin{remark}[Two root types and the Swiss-cheese coloring]
\label{rem:two-root-types}
The interior points $z_i$ carry the \emph{closed} color of the
Swiss-cheese operad $\SCchtop$; the boundary points $y_j^{(\ell)}$
carry the \emph{open} color.  The ordering on boundary points
is the $E_1$ (associative) structure; interior points are
unordered within each cluster (the $E_\infty$ or holomorphic
structure).  This two-coloring is the operadic fingerprint of the
$\mathrm{SC}^{\mathrm{ch},\mathrm{top}}$-algebra structure.
\end{remark}

\subsubsection{The bordered Fulton--MacPherson compactification}

\begin{construction}[Bordered FM compactification;
\ClaimStatusProvedHere]
\label{constr:bordered-FM}
\index{Fulton--MacPherson!bordered compactification|textbf}
\index{compactification!bordered FM|textbf}
The compactification
$\overline{\FM}^{\,\mathrm{bord}}_{g,r,s}(X)$ is obtained from
$X^r\times\prod_\ell(\partial_\ell X)^{s_\ell}$ by an iterated
sequence of real-oriented blowups, organized into three stages
reflecting the three collision regimes.

\medskip
\noindent\textbf{Stage I: Interior collisions.}
For every subset $S\subseteq\{1,\dotsc,r\}$ with $|S|\ge 2$,
blow up the diagonal
$\Delta_S=\{z_i=z_j\;\text{for all}\;i,j\in S\}$
in decreasing order of $|S|$ (total diagonal first, then
three-fold diagonals, etc.).  These are the standard FM blowups
of $\FM_r(X\setminus\partial X)$.  The exceptional divisor over
$\Delta_S$ is
\[
E_S^{\mathrm{int}}
\;\cong\;
\FM_{k_S}(X\setminus\partial X)\times\FM_{|S|}^{\mathrm{red}}(\mathbb C),
\]
where $k_S=r-|S|+1$ and $\FM_{|S|}^{\mathrm{red}}(\mathbb C)$ records
infinitesimal relative positions within the cluster, modulo
translation and dilation.  Each $E_S^{\mathrm{int}}$ is a
codimension-$1$ boundary divisor.

\medskip
\noindent\textbf{Stage II: Boundary collisions.}
For every boundary component $\partial_\ell X$ and every subset
$T\subseteq\{1,\dotsc,s_\ell\}$ with $|T|\ge 2$, blow up the
\emph{boundary diagonal}
\[
\Delta_T^{\partial,\ell}
\;=\;
\bigl\{y_i^{(\ell)}=y_j^{(\ell)}\;\text{for all}\;i,j\in T\bigr\}
\;\subseteq\;
(\partial_\ell X)^{s_\ell}.
\]
Since $\partial_\ell X$ is one-dimensional, the collision is a
real-codimension-$|T|-1$ locus and the blowup produces an
exceptional divisor
\[
E_T^{\partial,\ell}
\;\cong\;
\FM_{k_T}^{\partial}(\partial_\ell X)\times K_{|T|},
\]
where $K_{|T|}$ is the Stasheff associahedron of dimension $|T|-2$,
recording the order-preserving infinitesimal arrangement of $|T|$
points on a line.  This is the key distinction from the interior
case: the associahedron $K_m$ replaces the reduced FM space
$\FM_m^{\mathrm{red}}(\mathbb C)$, reflecting the passage from
holomorphic to topological (associative) factorization.

\medskip
\noindent\textbf{Stage III: Mixed collisions (interior approaching
boundary).}
For every pair $(i,j,\ell)$ with $i\in\{1,\dotsc,r\}$ and
$j\in\{1,\dotsc,s_\ell\}$, there is a codimension-$2$ locus
\[
\Delta_{i,j}^{\mathrm{mix},\ell}
\;=\;
\{z_i=y_j^{(\ell)}\in\partial_\ell X\}.
\]
More generally, for a subset $S\subseteq\{1,\dotsc,r\}$ and
$T\subseteq\{1,\dotsc,s_\ell\}$ with $|S|+|T|\ge 2$ and $|S|\ge 1$,
$|T|\ge 1$, the mixed diagonal is
\[
\Delta_{S,T}^{\mathrm{mix},\ell}
\;=\;
\bigl\{z_i=y_j^{(\ell)}\;\text{for all}\;i\in S,\;j\in T\bigr\}.
\]
Blow up these loci in the proper transform, ordered by decreasing
$|S|+|T|$.  The exceptional divisor is
\[
E_{S,T}^{\mathrm{mix},\ell}
\;\cong\;
\FM^{\mathrm{bord},\mathrm{red}}_{0,|S|,|T|}
\bigl(\mathbb H\bigr),
\]
where $\mathbb H=\{z\in\mathbb C:\operatorname{Im}(z)\ge 0\}$ is the
upper half-plane model, and the superscript ``red'' denotes the
quotient by translations along $\partial\mathbb H=\mathbb R$ and
dilations.  This records the infinitesimal arrangement of $|S|$
interior points and $|T|$ ordered boundary points within a
half-plane screen.

\medskip
\noindent\textbf{Result.}
The iterated blowup
$\overline{\FM}^{\,\mathrm{bord}}_{g,r,s}(X)$ is a compact smooth
manifold with corners.  Its codimension-$1$ boundary consists of the
divisors $E_S^{\mathrm{int}}$, $E_T^{\partial,\ell}$, and
$E_{S,T}^{\mathrm{mix},\ell}$ described above.
\end{construction}

\begin{remark}[Comparison with the Kontsevich--Soibelman and
Horel models]
\label{rem:KS-Horel-comparison}
\index{Swiss-cheese!bordered compactification}
When $X=\mathbb H$ (the closed upper half-plane),
$\overline{\FM}^{\,\mathrm{bord}}_{0,r,s}(\mathbb H)$ recovers the
Kontsevich--Soibelman \cite{kontsevich-soibelman} and Horel
\cite{horel-swiss-cheese} compactification of the two-colored
Swiss-cheese configuration space.  The three blowup stages
correspond to the three colors of the Voronov Swiss-cheese operad:
interior-interior (bulk), boundary-boundary (brane), and
interior-boundary (bulk-to-boundary).
\end{remark}

\subsubsection{Boundary strata and planted forests with two root types}

\begin{definition}[Two-rooted planted forests;
\ClaimStatusProvedHere]
\label{def:two-rooted-forest}
\index{planted forest!two-rooted|textbf}
A \emph{two-rooted planted forest} of type $(g,r,s)$ is a
rooted forest $\mathcal T$ whose vertices carry the data:
\begin{enumerate}[label=(\alph*)]
\item each vertex $v$ is labeled \emph{interior} ($\circ$-type) or
  \emph{boundary} ($\bullet$-type);
\item interior vertices have genus labels $g(v)\ge 0$ summing to
  $g$;
\item the leaves are partitioned into interior leaves (labeled by
  $\{1,\dotsc,r\}$) and boundary leaves (labeled by
  $\{1,\dotsc,s_\ell\}$ for each boundary component $\ell$);
\item \textbf{Swiss-cheese directionality constraint}:
  every edge from a $\bullet$-vertex to a $\circ$-vertex is
  \emph{forbidden}.  Equivalently, $\bullet$-vertices may only
  be descendants of $\circ$-vertices, never ancestors.
  This encodes the strict SC directionality: open-to-closed
  composition is EMPTY.
\item boundary vertices respect the linear ordering on each
  boundary component: the children of a $\bullet$-vertex inherit the
  planar order from $\partial_\ell X$.
\end{enumerate}
\end{definition}

\begin{proposition}[Boundary stratification;
\ClaimStatusProvedHere]
\label{prop:bordered-stratification}
\index{Fulton--MacPherson!bordered!boundary strata}
The codimension-$k$ boundary strata of
$\overline{\FM}^{\,\mathrm{bord}}_{g,r,s}(X)$ are in canonical
bijection with two-rooted planted forests $\mathcal T$ of type
$(g,r,s)$ with $k$ internal edges.  The bijection sends a
forest~$\mathcal T$ to the stratum
\[
\mathcal S_{\mathcal T}
\;\cong\;
\prod_{v\;\circ\text{-type}}
\FM_{n(v)}^{\mathrm{red}}(\mathbb C)
\;\times\;
\prod_{v\;\bullet\text{-type}}
K_{n(v)}
\;\times\;
\prod_{e\;\mathrm{mixed}}
\overline{\FM}^{\,\mathrm{bord},\mathrm{red}}_{0,n_\circ(e),n_\bullet(e)}
(\mathbb H),
\]
where $n(v)$ is the valence of vertex $v$, $K_{n(v)}$ is the
Stasheff associahedron, and the mixed product runs over edges
connecting $\circ$-type and $\bullet$-type vertices.
\end{proposition}

\begin{proof}
The proof is by direct inspection of the iterated blowup
construction.  Each Stage~I blowup introduces a boundary stratum
indexed by a subset $S$ of interior points forming a cluster,
corresponding to an internal edge of a $\circ$-rooted tree.
Each Stage~II blowup introduces a stratum indexed by a subset $T$
of ordered boundary points, corresponding to an internal edge of
a $\bullet$-rooted tree.  The associahedron $K_{|T|}$ arises
because boundary points are linearly ordered: the real blowup
of ordered coincidences on $\mathbb R$ produces precisely the
Stasheff polytope (this is the Axelrod--Singer--Fulton--MacPherson
identification in the real one-dimensional case;
cf.~Stasheff~\cite{stasheff-associahedra}).

Stage~III mixed blowups correspond to nested collisions
where interior points approach boundary points.  The exceptional
divisor records the relative positions in a half-plane screen
$\mathbb H$, which is the operadic composition space of the
Swiss-cheese operad.  The directionality constraint (no
$\bullet$-to-$\circ$ edges) holds because the blowup geometry of
a point on the boundary of $X$ does not produce an exceptional
divisor in the interior: the half-plane screen has the boundary on
its edge, so the resulting infinitesimal configuration is
a bulk-to-boundary composition, never the reverse.

Codimension-$k$ strata correspond to intersections of $k$ boundary
divisors, i.e.\ forests with $k$ internal edges.  The product
decomposition follows from the local product structure of the FM
blowup along each stratum.
\end{proof}

\subsubsection{Operadic structure}

\begin{theorem}[Bordered Swiss-cheese operadic structure;
\ClaimStatusProvedHere]
\label{thm:bordered-operad}
\index{Swiss-cheese operad!bordered operadic structure|textbf}
The collection
$\bigl\{\overline{\FM}^{\,\mathrm{bord}}_{0,r,s}(X)
\bigr\}_{r,s}$ carries a two-colored operadic structure with:
\begin{enumerate}[label=(\roman*)]
\item \textbf{Closed-color composition}:
  for subsets $S\subseteq\{1,\dotsc,r\}$ with $|S|=k$, the
  composition map
  \[
  \gamma_S^{\mathrm{int}}\colon
  \overline{\FM}^{\,\mathrm{bord}}_{0,r-k+1,s}(X)
  \times
  \FM_k^{\mathrm{red}}(\mathbb C)
  \longrightarrow
  \partial_{E_S^{\mathrm{int}}}
  \overline{\FM}^{\,\mathrm{bord}}_{0,r,s}(X)
  \]
  is the inclusion of the boundary divisor $E_S^{\mathrm{int}}$,
  and on chain complexes it is the standard FM composition
  governing the holomorphic factorization.

\item \textbf{Open-color composition}:
  for subsets $T\subseteq\{1,\dotsc,s_\ell\}$ with $|T|=m$, the
  composition map
  \[
  \gamma_T^{\partial}\colon
  \overline{\FM}^{\,\mathrm{bord}}_{0,r,s-m+1}(X)
  \times
  K_m
  \longrightarrow
  \partial_{E_T^{\partial,\ell}}
  \overline{\FM}^{\,\mathrm{bord}}_{0,r,s}(X)
  \]
  is the inclusion of the boundary divisor $E_T^{\partial,\ell}$,
  where $K_m$ is the Stasheff associahedron.  On chain complexes
  this governs the associative (topological, $E_1$) factorization.

\item \textbf{Mixed composition (Swiss-cheese insertion)}:
  for mixed subsets with $|S|\ge 1$ interior and $|T|\ge 1$
  boundary points, the composition
  \[
  \gamma_{S,T}^{\mathrm{mix}}\colon
  \overline{\FM}^{\,\mathrm{bord}}_{0,r',s'}(X)
  \times
  \overline{\FM}^{\,\mathrm{bord},\mathrm{red}}_{0,|S|,|T|}
  (\mathbb H)
  \longrightarrow
  \partial_{E_{S,T}^{\mathrm{mix},\ell}}
  \overline{\FM}^{\,\mathrm{bord}}_{0,r,s}(X)
  \]
  encodes the bulk-to-boundary coupling.

\item \textbf{Strict directionality}:
  There is no composition map with a $\bullet$-type (boundary)
  input producing a $\circ$-type (interior) output.  The open-to-%
  closed direction of the Swiss-cheese operad is empty.  This
  follows from the geometric fact that the exceptional divisor
  $E_{S,T}^{\mathrm{mix},\ell}$ lives over the boundary
  $\partial X$, not in the interior: there is no blowup
  mechanism for boundary clusters to migrate into the bulk.
\end{enumerate}
\end{theorem}

\begin{proof}
The proof reduces to checking that the boundary inclusions are
compatible with the operad axioms (equivariance, associativity of
compositions, unitality).

\emph{Closed-color associativity.}
The nested composition $\gamma_{S'}^{\mathrm{int}}\circ
(\id\times\gamma_{S''}^{\mathrm{int}})$ factors through a
codimension-$2$ stratum $E_{S'}\cap E_{S''}$ of the FM
compactification.  The two orderings of blowing up (first $S''$
inside $S'$, or first $S'$ then $S''$) produce canonically
isomorphic spaces by the standard FM argument (Axelrod--Singer
\cite{axelrod-singer}, Proposition~3.7), so the operad axiom
$\gamma\circ(\id\otimes\gamma)=\gamma\circ(\gamma\otimes\id)$
holds.

\emph{Open-color associativity.}
The nested boundary composition
$\gamma_{T'}^{\partial}\circ(\id\times\gamma_{T''}^{\partial})$
factors through the codimension-$2$ face of the associahedron.
The Stasheff associahedron $K_m$ is itself the operad of the
$A_\infty$-operad \cite{stasheff-associahedra}, and its face
identifications are the $A_\infty$ composition axioms.  Hence
open-color compositions inherit associativity from the
associahedron.

\emph{Mixed composition and directionality.}
A mixed composition followed by a closed-color composition is a
codimension-$2$ stratum of type
$E_{S}^{\mathrm{int}}\cap E_{S',T'}^{\mathrm{mix}}$.  Associativity
holds by the same iterated-blowup argument applied to the
half-plane model $\mathbb H$.  Conversely, a closed-color
composition cannot follow an open-color composition to produce
an interior point, because the half-plane screen
$\overline{\FM}^{\,\mathrm{bord},\mathrm{red}}_{0,|S|,|T|}(\mathbb H)$
has the interior strictly above $\partial\mathbb H$: there is no
continuous family of boundary configurations that limits to an
interior configuration.  This is the geometric source of the
strict Swiss-cheese directionality.
\end{proof}

\subsubsection{The diagonal bicomodule from bordered chains}

\begin{definition}[Bordered chain complex;
\ClaimStatusProvedHere]
\label{def:bordered-chains}
\index{diagonal bicomodule!geometric realization|textbf}
Define the \emph{bordered chain complex} in bidegree $(r,s)$
(with $s\in\mathbb Z_{\ge 0}$ corresponding to a single boundary
component, $b=1$):
\[
C_\Delta^{\mathrm{geom}}(r,s)
\;:=\;
C_*\bigl(
\overline{\FM}^{\,\mathrm{bord}}_{0,r,s}(X);\,
\omega^{\mathrm{bord}}_{r,s}
\bigr),
\]
where $\omega^{\mathrm{bord}}_{r,s}$ is the logarithmic weight
form defined as follows.  On the interior factor, use the standard
FM logarithmic forms $\omega^{\mathrm{hol}}$ (the weight forms
of the chiral algebra $\cA$, with logarithmic poles along
$E_S^{\mathrm{int}}$).  On the boundary factor, use the topological
weight forms $\omega^{\mathrm{top}}$ (which are piecewise constant
on each chamber of $K_m$, encoding the ordered product).  On the
mixed factor, use the Swiss-cheese weight form
$\omega^{\mathrm{SC}}$ (the bulk-to-boundary coupling form,
logarithmic along $E_{S,T}^{\mathrm{mix}}$ with residues given by
the $\beta$-map of Construction~\ref*{V1-constr:beta-map}).
\end{definition}

\begin{theorem}[Bicomodule structure on bordered chains;
\ClaimStatusProvedHere]
\label{thm:bordered-bicomodule}
\index{bicomodule!bordered geometric realization}
The collection $\{C_\Delta^{\mathrm{geom}}(r,s)\}_{r,s}$ carries a
canonical $C$-bicomodule structure, where
$C=\Barch(\cA)=\bigoplus_r C_*(
\overline{\FM}_{0,r,0}(X\setminus\partial X);\omega^{\mathrm{hol}}_r)$
is the interior bar coalgebra.  Specifically:

\begin{enumerate}[label=(\alph*)]
\item The \textbf{left coaction}
  (closed-color, holomorphic)
  \[
  \lambda\colon
  C_\Delta^{\mathrm{geom}}(r,s)
  \longrightarrow
  \bigoplus_{p+q=r}
  C(p)\otimes C_\Delta^{\mathrm{geom}}(q,s)
  \]
  is defined by the residue map along the interior boundary
  divisors $E_S^{\mathrm{int}}$: a configuration in
  $\overline{\FM}^{\,\mathrm{bord}}_{0,r,s}(X)$ approaching a
  Stage~I boundary stratum splits into a closed-color cluster
  (contributing to~$C$) and a residual bordered configuration
  (contributing to $C_\Delta^{\mathrm{geom}}$).

\item The \textbf{right coaction}
  (open-color, topological)
  \[
  \rho\colon
  C_\Delta^{\mathrm{geom}}(r,s)
  \longrightarrow
  \bigoplus_{p+q=s}
  C_\Delta^{\mathrm{geom}}(r,p)\otimes C^{\partial}(q)
  \]
  is defined by the residue map along the boundary divisors
  $E_T^{\partial}$, where
  $C^{\partial}(q)=C_*(K_q)$ is the Stasheff associahedron
  chain complex (the $E_1$-coalgebra of the boundary).
  A configuration approaching a Stage~II boundary stratum splits
  into a residual bordered configuration and a boundary cluster
  encoding the ordered product.

\item The left and right coactions commute:
  $(\lambda\otimes\id)\circ\rho
  =(\id\otimes\rho)\circ\lambda$.
  This holds because the Stage~I and Stage~II blowups are
  independent (interior and boundary diagonals are disjoint
  loci in $X^r\times(\partial X)^s$), so the two residue maps
  commute on the nose.
\end{enumerate}
\end{theorem}

\begin{proof}
\emph{Coassociativity of $\lambda$.}
The left coaction $\lambda$ is coassociative because the
iterated residue along $E_{S'}\cap E_{S''}$ factors as
$(\Delta_C\otimes\id)\circ\lambda$, by the standard FM
composition associativity of
Theorem~\ref{thm:bordered-operad}(i).  The counit is the
inclusion $C_\Delta^{\mathrm{geom}}(0,s)\hookrightarrow
C_\Delta^{\mathrm{geom}}(r,s)$ at $r=0$.

\emph{Coassociativity of $\rho$.}
The right coaction $\rho$ is coassociative because the iterated
residue along $E_{T'}\cap E_{T''}$ factors as
$(\id\otimes\Delta_{C^\partial})\circ\rho$, by the associahedron
face relations (Theorem~\ref{thm:bordered-operad}(ii)).
The counit is the inclusion at $s=0$.

\emph{Bicomodule compatibility.}
The commutativity $(\lambda\otimes\id)\circ\rho
=(\id\otimes\rho)\circ\lambda$ reduces to the independence of
the two blowup stages: the interior diagonal $\Delta_S$ and the
boundary diagonal $\Delta_T^\partial$ are disjoint closed
submanifolds of $X^r\times(\partial X)^s$ (because
$X\setminus\partial X$ and $\partial X$ are disjoint), so
the two blowups commute and the two residue operations commute
on the chain level.

\emph{Differential compatibility.}
The bordered chain complex carries a differential from the
boundary operator $\partial$ on chains plus the bar differential
$d_{\mathrm{bar}}$ from OPE residues along interior collision
divisors plus the $A_\infty$ differential $d_{A_\infty}$ from
boundary collision divisors.  The compatibility
$d^2=0$ follows from Stokes' theorem on the manifold with corners
$\overline{\FM}^{\,\mathrm{bord}}_{0,r,s}(X)$, exactly as in
Theorem~\ref{thm:Stokes_FM} of \S\ref{sec:FM_calculus}: the
sum of all codimension-$1$ residue contributions vanishes,
and codimension-$2$ corners cancel by the bordered analogue
of the Arnold--Orlik--Solomon relations.
\end{proof}

\subsubsection{Residue maps: OPE, product, and bulk-to-boundary}

\begin{construction}[Three residue maps;
\ClaimStatusProvedHere]
\label{constr:three-residues}
\index{residue map!bordered|textbf}
The bordered FM compactification carries three types of
codimension-$1$ residue maps, one for each collision regime.

\medskip
\noindent\textbf{(i) Interior residue (OPE, closed color).}
Along the divisor $E_S^{\mathrm{int}}$, the logarithmic weight
form $\omega^{\mathrm{bord}}_{r,s}$ has a log pole with
residue
\[
\Res_{E_S^{\mathrm{int}}}(\omega^{\mathrm{bord}}_{r,s})
\;=\;
r^{\mathrm{coll}}_S\bigl(\omega^{\mathrm{hol}}_{|S|}\bigr)
\;\otimes\;
\omega^{\mathrm{bord}}_{r-|S|+1,\,s},
\]
where $r^{\mathrm{coll}}_S$ is the collision residue extracting the
OPE coefficient, exactly as in the unbounded case
(Definition~\ref{def:residue_divisor} in \S\ref{sec:FM_calculus}).
This is the holomorphic bar differential: it encodes the chiral
product of the closed-color algebra.

\medskip
\noindent\textbf{(ii) Boundary residue (ordered product, open
color).}
Along the divisor $E_T^{\partial,\ell}$, the residue extracts
the ordered product of boundary fields:
\[
\Res_{E_T^{\partial,\ell}}(\omega^{\mathrm{bord}}_{r,s})
\;=\;
\mu_T^{\partial}\bigl(\omega^{\mathrm{top}}_{|T|}\bigr)
\;\otimes\;
\omega^{\mathrm{bord}}_{r,\,s-|T|+1},
\]
where $\mu_T^{\partial}$ is the $A_\infty$ multiplication map
on the boundary algebra $\cA^! = \cA|_{\partial X}$
(the Koszul dual, which lives on the boundary,
not the bulk).  The Stasheff associahedron face structure
ensures that the residues compose associatively.

\medskip
\noindent\textbf{(iii) Mixed residue (bulk-to-boundary
$\beta$-map).}
Along the divisor $E_{S,T}^{\mathrm{mix},\ell}$, the residue
encodes the bulk-to-boundary coupling:
\[
\Res_{E_{S,T}^{\mathrm{mix},\ell}}
(\omega^{\mathrm{bord}}_{r,s})
\;=\;
\beta_{S,T}\bigl(\omega^{\mathrm{SC}}_{|S|,|T|}\bigr)
\;\otimes\;
\omega^{\mathrm{bord}}_{r-|S|,\,s-|T|+1},
\]
where $\beta_{S,T}$ is the bulk-to-boundary restriction map:
it takes $|S|$ interior fields and $|T|$ boundary fields and
produces a single boundary output via the Swiss-cheese composition.
The half-plane screen $\mathbb H$ gives the local model: interior
points in $\{z\in\mathbb C:\operatorname{Im}(z)>0\}$ approach
boundary points on $\operatorname{Im}(z)=0$, and the residue
is the radial integral over the half-circle
$S^1_+\cap\mathbb H$.

\medskip
The strict SC directionality is manifest: there is no residue map
producing an interior output from boundary inputs, because
$E_{S,T}^{\mathrm{mix},\ell}$ lies over $\partial X$ and the
output point is on the boundary.
\end{construction}

\subsubsection{Identification with the algebraic diagonal bicomodule}

\begin{theorem}[Geometric realization of $C_\Delta$;
\ClaimStatusProvedHere]
\label{thm:bordered-realizes-diagonal}
\index{diagonal bicomodule!geometric realization!identification}
Let $\cA$ be a strongly admissible associative chiral algebra on $X$,
and set $C=\Barch(\cA)$.  Then there is a quasi-isomorphism of dg
$C$-bicomodules
\[
C_\Delta^{\mathrm{geom}}
\;\xrightarrow{\;\sim\;}
C_\Delta
\]
identifying the geometrically defined bordered bicomodule of
Definition~\textup{\ref{def:bordered-chains}} with the
algebraically defined diagonal bicomodule of the ordered bar
coalgebra.

More precisely, the following hold:
\begin{enumerate}[label=\textup{(\arabic*)}]
\item \textbf{Interval component.}
  For $X=[0,1]\times\mathbb C$ with a single boundary component,
  the bordered chain complex $C_\Delta^{\mathrm{geom}}$ reproduces
  $C_\Delta$ as a $C$-bicomodule.

\item \textbf{Gluing.}
  Gluing two bordered surfaces along a common boundary interval
  is computed by derived cotensor:
  \[
  C_\Delta^{\mathrm{geom}}(X_1)
  \square_C^{\mathbf R}
  C_\Delta^{\mathrm{geom}}(X_2)
  \;\simeq\;
  C_\Delta^{\mathrm{geom}}(X_1\cup_{\partial}X_2).
  \]

\item \textbf{Ordered pair-of-pants.}
  The ordered pair-of-pants cobordism $P$ (a disk with two
  incoming boundary intervals and one outgoing) induces
  \[
  \mu_P^{\mathrm{geom}}\colon
  C_\Delta^{\mathrm{geom}}\star C_\Delta^{\mathrm{geom}}
  \longrightarrow
  C_\Delta^{\mathrm{geom}},
  \]
  and under the identification with $C_\Delta$ this agrees with
  the algebraic pair-of-pants multiplication $\mu_P$ of
  Theorem~\textup{\ref{thm:pair-of-pants}}.

\item \textbf{Open trace formalism.}
  The full ordered genus-$0$ open trace formalism of
  Theorem~\textup{\ref{thm:ordered-open}} is recovered by the
  geometric operations on
  $\overline{\FM}^{\,\mathrm{bord}}_{0,r,s}(X)$.
\end{enumerate}
\end{theorem}

\begin{proof}
\emph{Statement~(1): interval identification.}
For $X=[0,1]\times\mathbb C$, the bordered configuration space
parametrizes $r$ interior points in
$(0,1)\times\mathbb C$ and $s$ ordered boundary points on
$\{0,1\}\times\mathbb C$.  The interior-to-boundary projection
$(t,z)\mapsto z$ gives an identification of the interior
configuration with
$\FM_r(\mathbb C)\times\Conf_r^{\mathrm{ord}}(\mathbb R)$
(the chiral-topological product), and the boundary points
contribute $\Conf_s^{\mathrm{ord}}(\mathbb R)$.
The resulting chain complex is
$C\widehat\otimes\cA\widehat\otimes C$ with the two-sided
twisting differential of Definition~\ref{def:Kbi}, because
the interior residues give the bar differential on the two
outer copies and the mixed residues give the twisting terms
$d_L$ and $d_R$.  By Theorem~\ref{thm:tangent=K} and
Theorem~\ref{thm:diagonal}, this is $\KK_\cA^{\mathrm{bi}}(\cA)
\simeq C_\Delta$.

\emph{Statement~(2): gluing.}
Gluing two bordered surfaces $X_1,X_2$ along a common boundary
interval identifies boundary points of $X_1$ with boundary
points of $X_2$.  On the chain level, this identification is
computed by the derived cotensor product over $C$ (the
coalgebraic analogue of the tensor product over $A$),
because the shared boundary carries the bar coalgebra $C$ and
the identification is the coalgebraic codiagonal.  This is
exactly the content of Corollary~\ref{cor:tensor-cotensor}.

\emph{Statement~(3): pair-of-pants.}
The pair-of-pants $P$ is a disk with three boundary segments.
The bordered FM compactification
$\overline{\FM}^{\,\mathrm{bord}}_{0,r,s_1+s_2+s_3}(P)$
has residue maps encoding the fusion of two incoming boundary
intervals into one outgoing.  The algebraic content is the
multiplication $\mu_A\colon A\widehat\otimes A\to A$, which
under $\KK_A^{\mathrm{bi}}$ gives
$\mu_P\colon C_\Delta\star C_\Delta\to C_\Delta$
(Theorem~\ref{thm:pair-of-pants}).  The geometric pair-of-pants
map is the composition of the gluing (Statement~(2)) with the
boundary contraction, which reproduces $\mu_P$ by the
uniqueness of the composition structure on the diagonal
(Corollary~\ref{cor:unit}).

\emph{Statement~(4): open trace formalism.}
The full open trace formalism consists of interval composition
(Statement~(2)), ordered pair-of-pants (Statement~(3)),
and the annulus action (Corollary~\ref{cor:cap}).  The first
two are proved above.  The annulus action is the geometric cap
product on bordered chains, transported to $C_\Delta$ by the
identification of Statement~(1); this reproduces
Corollary~\ref{cor:cap} by functoriality.  The sewing relations
are automatic because they hold on the algebraic side
(Theorem~\ref{thm:ordered-open}) and the identification is a
quasi-isomorphism of bicomodules.
\end{proof}

\subsubsection{Genus-$g$ extension and clutching}

\begin{construction}[Genus-$g$ bordered FM space;
\ClaimStatusProvedHere]
\label{constr:genus-g-bordered}
\index{Fulton--MacPherson!bordered!genus g@genus $g$|textbf}
For a closed surface $\Sigma_g$ of genus~$g$ with $b$ boundary
components, the bordered FM space
$\overline{\FM}^{\,\mathrm{bord}}_{g,r,s}(\Sigma_g)$ is constructed
by the same three-stage blowup procedure applied to $\Sigma_g$.
The new feature at genus $g\ge 1$ is the \emph{clutching map}:
cutting $\Sigma_g$ along a separating or non-separating simple
closed curve $\gamma$ produces either two surfaces of lower genus
or one surface of genus $g-1$ with two new boundary components.

\medskip
\noindent\textbf{Clutching maps.}
For a non-separating curve $\gamma\subset\Sigma_g$, cutting
produces $\Sigma_{g-1}$ with two new boundary circles
$\partial_\gamma^+,\partial_\gamma^-$.  The clutching map is
\[
\mu_\gamma\colon
\overline{\FM}^{\,\mathrm{bord}}_{g-1,r,s\cup(n,n)}(\Sigma_{g-1})
\longrightarrow
\overline{\FM}^{\,\mathrm{bord}}_{g,r,s}(\Sigma_g),
\]
obtained by identifying boundary points on $\partial_\gamma^+$ with
those on $\partial_\gamma^-$, summed over $n\ge 0$ with
$(-1)^n$ signs from the grading.  On the chain level this is
computed by derived cotensor over $C$:
\[
C_\Delta^{\mathrm{geom}}(\Sigma_g)
\;\simeq\;
C_\Delta^{\mathrm{geom}}(\Sigma_{g-1})
\square_C^{\mathbf R}
C_\Delta
\;\simeq\;
\coHoch^{\mathrm{ch}}_\bullet(C)\circ(\text{lower genus}).
\]
\end{construction}

\begin{theorem}[Genus-$g$ line-side modular $R$-matrix;
\ClaimStatusProvedHere]
\label{thm:genus-g-R-matrix-bordered}
\index{modular R-matrix!genus-g line-side|textbf}
\index{R-matrix!modular!from bordered FM|textbf}
From the bordered FM space
$\overline{\FM}^{\,\mathrm{bord}}_{g,2,0}(\Sigma_g)$ (genus~$g$
with $2$ interior points and no boundary points), the monodromy of
the ordered bar connection around the handle loops of $\Sigma_g$
produces the genus-$g$ line-side modular $R$-matrix
\[
R_g(z)
\;=\;
\operatorname{Hol}_{\gamma_1,\dotsc,\gamma_{2g}}
\bigl(
\nabla^{\mathrm{bar}}\big|_{\FM^{\mathrm{bord}}_{g,2,0}}
\bigr)
\;\in\;
\operatorname{End}(V\otimes V)[\![z^{-1}]\!].
\]
The genus expansion
\[
R^{\mathrm{mod}}(z;\hbar)
\;=\;
\sum_{g\ge 0}\hbar^{2g}\,R_g(z)
\]
agrees with the modular $R$-matrix of
Construction~\textup{\ref{constr:modular-ordered-bar}} to all
orders in~$\hbar$.

The clutching maps provide the recursion: $R_g$ is determined by
$R_{g-1}$ via the cotensor gluing along a non-separating curve,
with the correction term arising from the period matrix
$\tau_{ij}$ of $\Sigma_g$ through the genus-$g$ propagator.
\end{theorem}

\begin{proof}
The bar connection $\nabla^{\mathrm{bar}}$ on the interior
configuration space $\FM_2(\Sigma_g)$ is the flat connection
whose holonomy along generators of
$\pi_1(\Conf_2^{\mathrm{ord}}(\Sigma_g))$ gives the
$R$-matrix (Remark~\ref{rem:why-ordered}).  At genus~$0$
the fundamental group is the pure braid group $P_2\cong\mathbb Z$
and the holonomy is the genus-$0$ $R$-matrix $R_0(z)=R(z)$.

At genus~$g$, the fundamental group has additional generators
from the $2g$ handle loops $\gamma_1,\dotsc,\gamma_{2g}$ of
$\Sigma_g$.  The holonomy along these generators produces the
genus corrections $R_g(z)$ whose leading term involves the
period matrix:
for $\widehat{\mathfrak g}_k$, the genus-$1$ correction
$R_1(z)=(\kappa/(k+h^\vee)^2)\,\wp(z;\tau)\,\Omega+\cdots$
agrees with the formula in
Construction~\ref{constr:modular-ordered-bar}.

The bordered FM space
$\overline{\FM}^{\,\mathrm{bord}}_{g,2,0}(\Sigma_g)$ provides the
compactification needed to regularize the holonomy integral:
the logarithmic weight forms have at worst log poles along the
boundary divisors, and the residue along each divisor gives a
contribution to $R_g$ that is computed by the clutching recursion
of Construction~\ref{constr:genus-g-bordered}.  The clutching
formula shows that $R_g$ is determined by $R_{g-1}$ via
\[
R_g(z)
\;=\;
\sum_{a,b=1}^{2g}
\Bigl(\Omega^{-1}\Bigr)^{ab}\;
\Res_{\gamma_a}\Res_{\gamma_b}
\bigl(R_{g-1}(z)\bigr)
\;+\;
(\text{lower-genus terms}),
\]
where the double residue is along the pair of handle loops and the
matrix $(\Omega^{-1})^{ab}$ is the inverse of the intersection
form on $H_1(\Sigma_g;\mathbb Z)$.

Agreement with Construction~\ref{constr:modular-ordered-bar} to all
orders in $\hbar$ follows because both constructions compute the
same holonomy integral: the modular ordered bar complex of
Construction~\ref{constr:modular-ordered-bar} is the chain
complex of $\overline{\FM}^{\,\mathrm{bord}}_{g,2,0}(\Sigma_g)$
with coefficients in the weight forms, and the genus expansion
of the $R$-matrix is the genus expansion of the holonomy.
\end{proof}

\subsubsection{Annular closure: the remaining frontier}

\begin{theorem}[Annular closure computes coHochschild homology;
\ClaimStatusProvedHere]
\label{thm:bordered-annular}
\label{conj:bordered-annular}
\index{annular closure!bordered FM|textbf}
\index{coHochschild homology!annular realization|textbf}
Let $A$ be a strongly admissible augmented associative chiral algebra,
$C = \Barch(A)$ its ordered bar coalgebra, and
$C_\Delta$ the diagonal bicomodule.
Let $X = S^1 \times [0,1]$ be the standard annulus.
Then the bordered chain complex
\[
C_\bullet\!\bigl(
\overline{\FM}{}^{\,\mathrm{bord}}_{0,r,s}(S^1\times[0,1]);\,
\omega_{\mathrm{bar}}
\bigr)
\;\simeq\;
\coHoch^{\mathrm{ch}}_\bullet(C,C_\Delta).
\]
More precisely, there exists an annular bordered FM compactification
$\overline{\FM}{}^{\,\mathrm{ann}}_{r,s}$ whose chain complex,
weighted by bar forms, is quasi-isomorphic to the coHochschild complex
of $C$ with coefficients in the diagonal bicomodule $C_\Delta$.
\end{theorem}

\begin{proof}
The proof proceeds in five steps.

\medskip\noindent
\textbf{Step~1: Annular bordered FM compactification.}
The bordered FM compactification
$\overline{\FM}{}^{\,\mathrm{bord}}_{0,r,s}([0,1] \times \R)$
of Volume~I
(\S\ref*{V1-sec:bordered-fm})
parametrises $r$ interior points and $s$ boundary points on
the strip $[0,1] \times \R$, with the boundary decomposed into
a left component $\{0\} \times \R$ and a right component
$\{1\} \times \R$.  Boundary points carry the linear ordering
inherited from~$\R$.

Define the \emph{annular bordered FM space}
$\overline{\FM}{}^{\,\mathrm{ann}}_{r,s}$
as the homotopy quotient of
$\overline{\FM}{}^{\,\mathrm{bord}}_{0,r,s}([0,1] \times \R)$
by the identification $\{0\} \times \R \sim \{1\} \times \R$ that
glues the two boundary components into a single circle
$[0,1]/\{0 \sim 1\} = S^1$.  Concretely, this is the real-oriented
blow-up of the compactified configuration space of $r$ points in
the interior and $s$ points on the boundary of the annulus
$S^1 \times [0,1]$, where boundary points are \emph{cyclically}
ordered by the $S^1$ coordinate.

The key new feature relative to the strip is the
\emph{wrap-around collision}: a pair of consecutive boundary points
$(y_i, y_{i+1})$ on $S^1$ may collide by one point wrapping past the
gluing seam.  The exceptional divisor for this collision is the same
local model as for non-wrap-around boundary collisions (a copy of
$\overline{\FM}{}_2(\R) \cong [0,1]$), but its position in the
global combinatorics of strata is different: it connects the last
boundary point to the first, closing the linear sequence into a
cyclic one.

\medskip\noindent
\textbf{Step~2: Boundary strata and cyclic planted forests.}
The codimension-one boundary strata of
$\overline{\FM}{}^{\,\mathrm{ann}}_{r,s}$ are of the same four
types as for the bordered strip
(Proposition~\ref{prop:bordered-log-FM-strata}):
\begin{enumerate}[label=(\roman*)]
\item interior collisions (bar differential $d_{\mathrm{res}}$),
\item consecutive boundary collisions ($A_\infty$ faces),
\item mixed bubbling (bulk-to-boundary restriction),
\item wrap-around boundary collision (the new cyclic face).
\end{enumerate}
On the strip, boundary strata of type~(ii) are indexed by
\emph{planted forests}: ordered sequences of rooted trees, one
tree per consecutive block of boundary points.  Passing from the
strip to the annulus replaces the linear ordering of trees by a
\emph{cyclic} ordering: the last tree's rightmost leaf is followed
cyclically by the first tree's leftmost leaf.  Thus the boundary
strata are indexed by \emph{necklaces of planted trees}
(cyclic planted forests).

\medskip\noindent
\textbf{Step~3: Chain complex as cyclic cobar complex.}
Fix the open colour.  On the strip, the chain complex weighted by
bar forms is (by items~(1)--(2) of
Theorem~\ref{thm:bordered-geometric} and the Vol~I construction)
\[
C_\bullet\!\bigl(
\overline{\FM}{}^{\,\mathrm{bord}}_{0,0,s}([0,1]\times\R);\,
\omega_{\mathrm{bar}}
\bigr)
\;\cong\;
C_\Delta \otimes C^{\otimes (s-1)},
\]
with the strip differential encoding the two-sided coaction of $C$
on $C_\Delta$: the left coaction $\Delta_L$ from boundary collisions
near $\{0\}\times\R$ and the right coaction $\Delta_R$ from boundary
collisions near $\{1\}\times\R$.  This is the two-sided cobar
complex computing $\Cotor^{C^e}(C_\Delta, -)$.

Upon annular identification, the $s$ boundary points become
cyclically ordered.  Decompose an annular configuration of $s$
boundary points on $S^1$ by cutting at a marked basepoint to
obtain a linear sequence of $s$ points on $[0,1)$.  The resulting
chain group at cyclic tensor degree $n$ is
\[
\mathrm{cyc}^n(C, C_\Delta)
\;:=\;
C_\Delta \widehat\otimes_{C} C^{\widehat\otimes\, n},
\]
where the cotensor $\widehat\otimes_C$ identifies the right coaction
of $C_\Delta$ with the left coaction on the first tensor factor of
$C^{\widehat\otimes\, n}$, and the left coaction of $C_\Delta$ with
the right coaction on the last tensor factor.  This
identification is precisely the cyclic gluing:
\begin{itemize}
\item the \emph{right} coaction of $C_\Delta$ (boundary collisions
  approaching the seam from the right in the strip picture) is
  glued to the comultiplication acting on the first $C$-factor
  (boundary collisions just past the seam);
\item the \emph{left} coaction of $C_\Delta$ (boundary collisions
  approaching the seam from the left) is glued to the
  comultiplication on the last $C$-factor (boundary collisions
  just before the seam).
\end{itemize}

The differential on $\mathrm{cyc}^\bullet(C, C_\Delta)$ has three
terms:
\begin{enumerate}[label=(\alph*)]
\item the \emph{internal differential} $d_{\mathrm{int}}$:
  the sum of the internal differentials of $C_\Delta$ and each
  $C$-factor (coming from interior collisions, type~(i));
\item the \emph{comultiplication differential} $d_\Delta$:
  for each $i = 1, \ldots, n$, insert a comultiplication
  $\Delta\colon C \to C \otimes C$ on the $i$th factor,
  increasing the cyclic tensor degree by one (coming from
  consecutive non-wrap-around boundary collisions, type~(ii));
\item the \emph{wrap-around differential} $d_{\mathrm{wrap}}$:
  the comultiplication applied at the gluing seam, inserting a new
  $C$-factor between the last and first positions in the cyclic
  word (coming from wrap-around collisions, type~(iv)).
\end{enumerate}
The total differential is
$d = d_{\mathrm{int}} + d_\Delta + d_{\mathrm{wrap}}$,
and the identity $d^2 = 0$ follows from $\partial^2 = 0$
on the manifold-with-corners
$\overline{\FM}{}^{\,\mathrm{ann}}_{r,s}$,
just as in the strip case.

\medskip\noindent
\textbf{Step~4: Identification with coHochschild homology.}
The coHochschild complex of a conilpotent dg coalgebra $C$ with
coefficients in a $C$-bicomodule $M$ is
(Definition~\ref*{V1-def:coHochschild},
cf.\ Theorem~\ref{thm:HH-coHH-homology})
\[
\coHoch^{\mathrm{ch}}_\bullet(C, M)
\;=\;
\Tot\Bigl(
M \widehat\otimes_{C^e} \overline{C}^{\widehat\otimes\, \bullet}
\Bigr),
\]
where $\overline C = \ker(\eps\colon C \to k)$ is the coaugmentation
coideal, the cotensor over $C^e = C \widehat\otimes C^{\cop}$
identifies the left $C$-coaction on $M$ with the right coaction on
the last $\overline C$-factor and the right $C$-coaction on $M$
with the left coaction on the first $\overline C$-factor, and the
cosimplicial structure maps are given by comultiplications.

We claim the cyclic cobar complex $\mathrm{cyc}^\bullet(C, C_\Delta)$
of Step~3 is canonically isomorphic to
$\coHoch^{\mathrm{ch}}_\bullet(C, C_\Delta)$.
At the level of graded modules:
\[
\mathrm{cyc}^n(C, C_\Delta)
\;=\;
C_\Delta \widehat\otimes_C C^{\widehat\otimes\, n}
\;\cong\;
C_\Delta \widehat\otimes_{C^e} \overline{C}^{\widehat\otimes\, n},
\]
where the isomorphism uses the splitting
$C = k \oplus \overline C$ from the coaugmentation and the
conilpotency of $C$ (which ensures the completed cotensor converges).
This isomorphism is an identity of underlying graded modules: the
cyclic cotensor over a single copy of $C$ (identifying right of
$C_\Delta$ with left of the first bar factor, and left of $C_\Delta$
with right of the last bar factor) equals the cotensor over
$C^e$ (which simultaneously identifies both coactions).

Under this isomorphism, the three components of the geometric
differential correspond to the three components of the coHochschild
differential:
\begin{itemize}
\item $d_{\mathrm{int}}$ maps to the internal differential of
  the coHochschild complex;
\item $d_\Delta$ maps to the cosimplicial coface maps
  $\delta_i\colon \overline{C}^{\widehat\otimes\, n}
  \to \overline{C}^{\widehat\otimes\, (n+1)}$
  for $i = 1, \ldots, n-1$ (interior comultiplications);
\item $d_{\mathrm{wrap}}$ maps to the two ``end'' coface maps
  $\delta_0$ and $\delta_n$ that involve the bicomodule coactions
  of $C_\Delta$.
\end{itemize}
The wrap-around differential, the
geometric novelty of the annular construction, encodes
the bicomodule coactions of $C_\Delta$ at the ``seam'' of the
cyclic word.  In the strip picture, these are the left and right
coactions, which act independently; the annular identification glues
them into the cyclic coface maps $\delta_0$ and $\delta_n$ that
close the cosimplicial structure.  This is the geometric
manifestation of the formula
\[
\coHoch^{\mathrm{ch}}_\bullet(C, C_\Delta)
\;=\;
C_\Delta \square_{C^e}^{\mathbf R} C_\Delta,
\]
already established algebraically in
Corollary~\ref{cor:annulus}.

Since the differentials match, we obtain a canonical isomorphism of
chain complexes
\[
C_\bullet\!\bigl(
\overline{\FM}{}^{\,\mathrm{ann}}_{0,s};\,
\omega_{\mathrm{bar}}
\bigr)
\;\cong\;
\coHoch^{\mathrm{ch}}_\bullet(C, C_\Delta).
\]
For configurations with $r > 0$ interior points, the interior
collisions contribute the bar differential in the holomorphic
direction, which acts on each $C$-factor and on $C_\Delta$.  This
is compatible with the coHochschild structure because the bar
differential commutes with the coalgebra structure
(it is a coderivation), so
\[
C_\bullet\!\bigl(
\overline{\FM}{}^{\,\mathrm{ann}}_{r,s};\,
\omega_{\mathrm{bar}}
\bigr)
\;\simeq\;
\coHoch^{\mathrm{ch}}_\bullet(C, C_\Delta).
\]

\medskip\noindent
\textbf{Step~5: Recovery of the open trace formalism at genus~$0$.}
We verify that the annular construction is consistent with the
algebraic open trace formalism of
Theorem~\ref{thm:ordered-open}.

The open trace formalism provides an algebraic annular closure
$\mathcal H_C^{\mathrm{cl}} \simeq \coHoch^{\mathrm{ch}}_\bullet(C)$
(Corollary~\ref{cor:annulus}), together with a cap action on
$C_\Delta$ (Corollary~\ref{cor:cap}).  By the identification of
Step~4, the geometric annular closure
$C_\bullet(\overline{\FM}{}^{\,\mathrm{ann}}_{r,s};\omega_{\mathrm{bar}})$
computes the same coHochschild complex, so it recovers the closed
object $\mathcal H_C^{\mathrm{cl}}$.

The cap action has a geometric interpretation: a configuration
on the annulus with one marked ``output'' boundary point is an
element of
$\coHoch^{\mathrm{ch}}_\bullet(C) \square_C^{\mathbf R} C_\Delta$,
and the geometric evaluation at the output point produces the
algebraic cap map.  This geometric description specialises the
interval composition $\square_C^{\mathbf R}$ of the open trace
formalism.

Finally, the pair-of-pants operation $\mu_P$ is recovered from the
annular construction by cutting the annulus along a boundary
arc to produce two half-annuli (= strips), glued along a shared
boundary interval.  This cutting is the geometric inverse of
the annular closure, and the resulting operation is the ordered
fusion product $\star$ of Definition~\ref{def:Kbi}
followed by interval composition, which is $\mu_P$
by Theorem~\ref{thm:pair-of-pants}.

Thus all genus-zero operations of the open trace formalism are
recovered geometrically.
\end{proof}

\begin{remark}[Proof strategy and the role of cyclic gluing]
\label{rem:bordered-annular-strategy}
The heart of the proof is the observation that the cyclic
identification $\{0\} \sim \{1\}$ on the strip intertwines the
left and right coactions of the diagonal bicomodule $C_\Delta$
in precisely the way prescribed by the coHochschild complex.
The wrap-around collision stratum (the only stratum without an
analogue in the strip) produces the ``closing'' coface maps
$\delta_0$ and $\delta_n$ that distinguish coHochschild from the
two-sided cobar complex.  Without wrap-around, the chain complex
would compute $\Cotor^{C^e}(C_\Delta, C_\Delta)$ (the derived
cotensor of two \emph{independent} copies of $C_\Delta$ over
$C^e$); with wrap-around, the two copies are identified and the
result is the self-cotrace, i.e.\ coHochschild homology.
\end{remark}

\begin{theorem}[Annular FM as genus-$1$ modular Swiss-cheese;
\ClaimStatusProvedHere]
\label{thm:bordered-annular-modular}
\index{annular FM!modular Swiss-cheese|textbf}
\index{modular operad!genus-1 cell|textbf}
The annular FM compactification
$\overline{\FM}{}^{\,\mathrm{ann}}_{r,s}$ is the genus-$1$ part
of the modular Swiss-cheese operad $\cM_{\SCchtop}$, in the
following precise sense:
\begin{enumerate}[label=\textup{(\roman*)}]
\item The unique genus-$1$ surface with one boundary component
is the annulus $S^1 \times [0,1]$.
\item The annular FM compactification parametrises configurations
on this surface with the correct collision/wrapping boundary
strata (Proposition~\ref{prop:bordered-log-FM-strata} plus
the wrapping faces of type~(iv)).
\item The composition maps of the modular operad (clutching along
boundary curves) are computed by derived cotensor over $C$,
matching the annular gluing.
\end{enumerate}

The clutching recursion connects genus~$0$ and genus~$1$: cutting
the annulus along a boundary arc produces the strip, and the
annular closure is the clutching of the two strip endpoints.
On chain complexes:
\[
B^{\mathrm{ann}}_\bullet(\cA)
\;\simeq\;
B^{\mathrm{strip}}_\bullet(\cA)\;\otimes_{C}\;k_{\mathrm{trace}},
\]
where $k_{\mathrm{trace}}$ is the ground field viewed as a
$C$-comodule via the coaugmentation (the cotrace map).
\end{theorem}

\begin{proof}
The annular FM space is defined by the identification
$\{0\} \sim \{1\}$ on the bordered strip FM space
(Definition~\ref{def:bordered-log-FM}).  This identification
is exactly the self-gluing operation of the modular operad
$\cM_{\SCchtop}$ applied to the genus-$0$ strip with its two
boundary components.  The resulting genus is $0 + 1 = 1$
(one handle created by the identification), and the boundary
count decreases by $2$ to $1$ (the two strip boundaries are
identified into a single circle $S^1$).  The strata of the
resulting space are computed by the boundary formula of
Proposition~\ref{prop:bordered-log-FM-strata}, with the
wrapping faces arising from the identification of the two
boundary strata of the strip.  The chain-level statement follows
from the identification of Theorem~\ref{thm:bordered-annular}:
the annular complex is the coHochschild complex, which is
$C_\Delta \square_{C^e}^{\mathbf R} C_\Delta$, and the
cotrace map is the coaugmentation $C \to k$ applied to
the strip complex.
\end{proof}

\begin{remark}[The annulus as self-cobordism]
\label{rem:annulus-self-cobordism}
The annulus $S^1 \times [0,1]$ is the identity cobordism of $S^1$
in the category of $1$-dimensional cobordisms.  The strip
$[0,1] \times \R$ is the identity cobordism of the interval $[0,1]$
in the category of $0$-dimensional cobordisms.  The passage from
strip to annulus is the passage from the $0{+}1$-dimensional open
TCFT (the ordered $A_\infty$ theory) to the $1{+}1$-dimensional open
TCFT (the Hochschild/cyclic theory).  The annular FM
compactification is the first nontrivial modular cell: it is the
genus-$1$ operation of the modular extension of the ordered
Swiss-cheese.
\end{remark}

\subsubsection{Genus-$1$ $R$-matrix from annular monodromy}

The annular FM compactification provides the genus-$1$ correction
to the $R$-matrix.  On the strip, the $R$-matrix
$R_0(z)$ is the monodromy of the ordered bar connection around
$\pi_1(\Conf_2^{\mathrm{ord}}(\C)) = \Z$.  On the annulus, the
fundamental group acquires an additional generator from the
$S^1$ factor:
\[
\pi_1(\Conf_2^{\mathrm{ord}}(S^1 \times \C))
\;=\;
P_2 \ltimes \pi_1(S^1)^2
\;\supset\;
P_2 = \pi_1(\Conf_2^{\mathrm{ord}}(\C)),
\]
where $P_2 \cong \Z$ is the pure braid group and the two copies
of $\pi_1(S^1) = \Z$ correspond to each point winding around
the $S^1$ factor.

\begin{theorem}[Genus-$1$ $R$-matrix from annular monodromy;
\ClaimStatusProvedHere]
\label{thm:genus1-R-matrix}
\index{R-matrix!genus-1 correction|textbf}
\index{Weierstrass function!genus-1 propagator|textbf}
The genus-$1$ correction to the $R$-matrix is the monodromy of
the ordered bar connection along the new generator:
\[
R_1(z;\tau)
\;=\;
\operatorname{Mon}_{\gamma_{S^1}}
\bigl(\nabla^{\mathrm{bar}}\big|_{\mathrm{ann}}\bigr),
\]
where $\gamma_{S^1}$ is the loop in which one point winds around
the $S^1$ factor.  For $\cA = \widehat{\fg}_k$:
\[
R_1(z;\tau)
\;=\;
\frac{\kappa(\widehat{\fg}_k)}{(k+h^\vee)^2}\,
\wp(z;\tau)\,\Omega
\;+\;
O(\wp'),
\]
where $\wp(z;\tau)$ is the Weierstrass $\wp$-function,
$\Omega = \sum_a e_a \otimes e^a$ is the Casimir element
of~$\fg$, and $\kappa(\widehat{\fg}_k) = k/(k+h^\vee)$ is the
curvature.  The full modular $R$-matrix is
\[
R^{\mathrm{mod}}(z;\tau,\hbar)
\;=\;
R_0(z) + \hbar^2\, R_1(z;\tau) + O(\hbar^4),
\]
in agreement with the clutching recursion of
Construction~\ref{constr:modular-ordered-bar}.
\end{theorem}

\begin{proof}
The genus-$1$ $R$-matrix arises from the monodromy of
the ordered bar connection $\nabla^{\mathrm{bar}}$ on the
annular configuration space.  The connection restricts from the
strip to the annulus by the identification of
Theorem~\ref{thm:bordered-annular-modular}.  For
$\cA = \widehat{\fg}_k$ at non-critical level, the connection
is the KZ connection
$\nabla_{\mathrm{KZ}} = d - \Omega\, d\log(z_1 - z_2)/(k + h^\vee)$
on the strip, which on the annulus becomes
$\nabla_{\mathrm{KZ}}^{\mathrm{ell}} = d - \Omega\, \varpi(z;\tau)/(k + h^\vee)$,
where $\varpi(z;\tau) = \partial_z \log \theta_1(z;\tau)$
is the elliptic propagator.  The monodromy around
$\gamma_{S^1}$ is computed by the Drinfeld--Kohno theorem
for elliptic KZ, giving
$R_1(z;\tau) = \kappa(\widehat{\fg}_k)\, \wp(z;\tau)\, \Omega / (k + h^\vee)^2$
at leading order, where $\wp = -\partial_z \varpi + \varpi^2$
is the Weierstrass function.  The agreement with
Construction~\ref{constr:modular-ordered-bar} follows because
both compute the same holonomy integral on the genus-$1$ surface.
\end{proof}

\begin{remark}[The Weierstrass $\wp$-function as genus-$1$ propagator]
\label{rem:wp-genus1-propagator}
The appearance of $\wp(z;\tau)$ is not coincidental: $\wp$ is the
genus-$1$ Green's function on $\C/(\Z \oplus \tau\Z)$, just as
$1/(z_1 - z_2)$ is the genus-$0$ Green's function on~$\C$.  The
passage from $1/z$ to $\wp(z;\tau)$ is the passage from the
genus-$0$ Cauchy kernel to the genus-$1$ Weierstrass kernel.
The modular properties of $\wp$ (quasi-periodicity under
$\tau \to \tau + 1$ and $\tau \to -1/\tau$) transmit to the modular
properties of $R_1$, which is the mathematical origin of modular
covariance of the genus-$1$ $R$-matrix.

This connects to the central open question of the genus-$1$
$E_1$ frontier: whether the curvature $\kappa(\cA) \cdot \omega_1$
and the braiding $R_1(z;\tau)$ decouple (as they do at genus~$0$)
or entangle (as the pole structure of $\wp$ suggests).  The
annular FM compactification provides the geometric arena where
this question is decided: the answer is read from the boundary
strata of $\overline{\FM}{}^{\,\mathrm{ann}}_{r,s}$ at
codimension~$2$, where the curvature divisor and the braiding
divisor may or may not intersect transversally.
\end{remark}

\subsection{Associative modular Maurer--Cartan theory}

The higher-genus modular story admits an associative analogue.
The ``ordered cyclic compactification'' sought in the original
conjecture is the Feynman transform $F\!\Ass$ of the ribbon modular
operad (Definition~\ref*{V1-def:ribbon-modular-operad}, Vol~I).

\begin{theorem}[Associative modular Maurer--Cartan class;
\ClaimStatusProvedHere]
\label{conj:modular}
The $E_1$ modular convolution dg~Lie algebra
\[
{\gAmod}^{E_1}
\;:=\;
\prod_{\substack{g,n \\ 2g-2+n>0}}
\operatorname{Hom}\bigl(
\cM_{\Ass}(g,n),\,\operatorname{End}_\cA(n)
\bigr)
\]
\textup{(}Definition~\textup{\ref*{V1-def:e1-modular-convolution}},
Vol~I\textup{)}
carries a canonical Maurer--Cartan element
\[
\Theta_\cA^{E_1}
\;:=\;
D_\cA^{E_1} - d_0
\;\in\;
\MC\bigl({\gAmod}^{E_1}\bigr),
\]
whose linear term is the genus anomaly, whose arity-$2$ component is
the modular $R$-matrix $R^{\mathrm{mod}}(z;\hbar)$, and whose full
Maurer--Cartan equation expresses compatibility of all ordered
higher-genus boundary gluings.

Under $\Sigma_n$-averaging\textup{,}
$\operatorname{av}(\Theta_\cA^{E_1}) = \Theta_\cA$
recovers the symmetric MC element of
Theorem~\textup{\ref*{V1-thm:mc2-bar-intrinsic}}.
\end{theorem}

\begin{proof}
The MC property $D\Theta^{E_1} + \tfrac12[\Theta^{E_1},\Theta^{E_1}]
= 0$ is automatic from $D_{F\!\Ass}^2 = 0$
(Theorem~\ref*{V1-thm:fass-d-squared-zero}, Vol~I), by the bar-intrinsic
construction: $\Theta^{E_1} = D^{E_1}_\cA - d_0$ satisfies MC
because $(D^{E_1}_\cA)^2 = 0$.  This is the $E_1$ analogue of
Theorem~\ref*{V1-thm:mc2-bar-intrinsic}.
See Remark~\ref*{V1-rem:conj-modular-resolved} (Vol~I) and
Proposition~\ref{prop:R-canonical-vol2} for the all-genus extension
of the $E_1$ package.
\end{proof}

\subsection{BRST and Drinfeld--Sokolov compatibility}

The last frontier is compatibility between associative self-duality and reduction.

\begin{theorem}[Reduction commutes with associative chiral duality
\textup{(}principal case\textup{)}; \ClaimStatusProvedHere]
\label{conj:DS}
Let $Q_{DS}$ denote the principal Drinfeld--Sokolov reduction functor
on the category of strongly admissible associative chiral algebras
at non-critical level.  Assume:
\begin{enumerate}[label=(\alph*)]
\item $Q_{DS}$ is exact on the relevant filtered complexes;
\item $Q_{DS}$ preserves PBW filtrations;
\item the BRST differential $d_{\mathrm{BRST}}$ lifts to the ordered
  bar coalgebra.
\end{enumerate}
Then
\[
Q_{DS}(\cA^{!,E_1})\simeq Q_{DS}(\cA)^{!,E_1}.
\]
For every $\cA$-bimodule $M$\textup{:}
\[
Q_{DS}\bigl(\KK_\cA^{\mathrm{bi}}(M)\bigr)\simeq
\KK_{Q_{DS}(\cA)}^{\mathrm{bi}}\bigl(Q_{DS}(M)\bigr).
\]
\end{theorem}

\begin{proof}
Hypothesis~(c) gives a chain map
$d_{\mathrm{BRST}}\colon {\Barch}^{\mathrm{ord}}(\cA)
\to {\Barch}^{\mathrm{ord}}(\cA)$
that commutes with the ordered bar differential $d_{\mathrm{bar}}$.
The key identity is
$[d_{\mathrm{BRST}}, d_{\mathrm{bar}}] = 0$:
this holds because the BRST differential is an inner derivation of
the chiral algebra (it preserves the OPE), and the bar differential
is built from the OPE.  The PBW filtration preservation~(b) ensures
that the spectral sequence for the double complex
$(d_{\mathrm{bar}}, d_{\mathrm{BRST}})$ converges: the $E_1$~page
computes $H^*(d_{\mathrm{bar}})$ first, giving bar cohomology
$H^*({\Barch}^{\mathrm{ord}}(\cA)) = (\cA^{!,E_1})^\vee$; then the
$E_2$~page applies $H^*(d_{\mathrm{BRST}})$, giving
$Q_{DS}((\cA^{!,E_1})^\vee)$.

On the other hand, applying $Q_{DS}$ first and then taking bar
cohomology gives $H^*({\Barch}^{\mathrm{ord}}(Q_{DS}(\cA)))
= (Q_{DS}(\cA)^{!,E_1})^\vee$.  Exactness~(a) gives
a spectral-sequence comparison, and the convergence (from PBW
filtration preservation) gives the quasi-isomorphism.

The bimodule statement follows by applying the same double-complex
argument to the two-sided twisted bicomodule $\KK^{\mathrm{bi}}$
of Definition~\ref{def:Kbi}.

\emph{Principal-specific input.}
For principal DS reduction, the ghost BRST complex has
level-independent conformal weights, so the PBW filtration
compatibility~(b) is automatic at non-critical level (this is the
content of the Feigin--Frenkel theorem).  For non-principal
reduction, hypothesis~(b) becomes the nontrivial requirement;
this remains the frontier
(Conjecture~\ref*{V1-conj:ds-kd-arbitrary-nilpotent}, i.e.\ the
non-principal nilpotent corridor).
\end{proof}

\begin{remark}[Scope of the theorem]
\label{rem:ds-scope}
For principal DS reduction, all three hypotheses are verified
at non-critical level, giving an unconditional theorem for the
principal corridor.  For hook-type in type~$A$
(Theorem~\ref*{V1-thm:hook-transport-corridor}), hypothesis~(b) holds
by the Fehily--Creutzig--Linshaw--Nakatsuka--Sato filtration analysis.
The full non-principal case
remains conjectural
(Conjecture~\ref*{V1-conj:ds-kd-arbitrary-nilpotent}).
\end{remark}
\subsection{The subregular $W_k(\mathfrak{sl}_3,f_{\mathrm{sub}})$:
a non-principal ordered bar complex}
\label{subsec:ordered-bar-subregular}

\begin{computation}[The ordered bar complex of the
subregular $W$-algebra $W_k(\mathfrak{sl}_3,f_{\mathrm{sub}})$;
\ClaimStatusProvedHere]
\label{comp:ordered-bar-n2-sca}
\index{W-algebra@$\mathcal W$-algebra!subregular!ordered bar|textbf}
\index{N=2 superconformal!ordered bar complex|textbf}
\index{bar complex!ordered!subregular W-algebra@subregular $\mathcal W$-algebra|textbf}

\providecommand{\Gp}{G^+}
\providecommand{\Gm}{G^-}
\providecommand{\wsl}{\widehat{\mathfrak{sl}}}

This is the first complete computation of the ordered bar
complex for a \emph{non-principal} $W$-algebra.  The subregular
$W_k(\mathfrak{sl}_3,f_{\mathrm{sub}})$ is isomorphic to the
$\mathcal N=2$ superconformal algebra (SCA).  Three phenomena
appear here that have no analogue in the principal case:
\begin{enumerate}[label=(\roman*)]
\item the ghost central charge is \emph{level-dependent},
not constant;
\item the bar complex carries a $\mathbb Z$-grading from the
$U(1)$ charge $J$, decomposing the differential into
charge sectors;
\item the Koszul dual is a \emph{different} $W$-algebra
(associated to the transpose nilpotent), not the same
algebra at dual level.
\end{enumerate}

\medskip
\noindent\textbf{Part I: Generator data.}
The algebra $W_k(\mathfrak{sl}_3,f_{\mathrm{sub}})$ has
four strong generators:
\begin{center}
\renewcommand{\arraystretch}{1.15}
\begin{tabular}{lcccc}
\textbf{Field} & \textbf{Weight $h$} & \textbf{$U(1)$ charge $q$}
& \textbf{Statistics} & \textbf{Bar degree $|s^{-1}\Phi|$} \\
\hline
$J$ & $1$ & $0$ & bosonic & $0$ \\
$\Gp$ & $3/2$ & $+1$ & fermionic & $1$ \\
$\Gm$ & $3/2$ & $-1$ & fermionic & $1$ \\
$T$ & $2$ & $0$ & bosonic & $0$
\end{tabular}
\end{center}
The central charge is $c = 3(2k+1)/(k+3)$, and we write
$\ell = c/3 = (2k+1)/(k+3)$ for the normalised level.

The bar degree of a desuspended generator $s^{-1}\Phi$ is
$|s^{-1}\Phi| = |\Phi| - 1$, where $|\Phi|$ is the
\emph{cohomological} (statistics) degree: $|\Phi|=0$ for
bosonic and $|\Phi|=1$ for fermionic fields.  Thus
$|s^{-1}J| = |s^{-1}T| = -1$ (odd) and
$|s^{-1}\Gp| = |s^{-1}\Gm| = 0$ (even).

\medskip
\noindent\textbf{Part II: The complete OPE.}
The ten independent OPE pairs, written in standard form
$\Phi_i(z)\,\Phi_j(w) \sim \sum_n r_n^{ij}/(z{-}w)^{n+1}$:

\smallskip
\noindent\emph{(i) Virasoro subalgebra.}
\begin{equation}\label{eq:n2-TT}
T(z)\,T(w)\;\sim\;\frac{c/2}{(z{-}w)^4}
+\frac{2T(w)}{(z{-}w)^2}+\frac{\partial T(w)}{z{-}w}.
\end{equation}
\emph{(ii) Primary fields under Virasoro.}
\begin{align}
T(z)\,\Gp(w)&\;\sim\;\frac{\tfrac32\,\Gp(w)}{(z{-}w)^2}
+\frac{\partial \Gp(w)}{z{-}w},\label{eq:n2-TGp}\\
T(z)\,\Gm(w)&\;\sim\;\frac{\tfrac32\,\Gm(w)}{(z{-}w)^2}
+\frac{\partial \Gm(w)}{z{-}w},\label{eq:n2-TGm}\\
T(z)\,J(w)&\;\sim\;\frac{J(w)}{(z{-}w)^2}
+\frac{\partial J(w)}{z{-}w}.\label{eq:n2-TJ}
\end{align}
\emph{(iii) $U(1)$ current algebra.}
\begin{equation}\label{eq:n2-JJ}
J(z)\,J(w)\;\sim\;\frac{\ell}{(z{-}w)^2}.
\end{equation}
Note: $J(z)\,J(w)$ has a second-order pole with coefficient
$\ell = c/3$, reflecting the $U(1)$ anomaly.  The zeroth
product $J_{(0)}J = 0$: the $U(1)$ current is abelian.

\smallskip
\noindent\emph{(iv) $U(1)$ charges.}
\begin{align}
J(z)\,\Gp(w)&\;\sim\;\frac{\Gp(w)}{z{-}w},\label{eq:n2-JGp}\\
J(z)\,\Gm(w)&\;\sim\;\frac{-\Gm(w)}{z{-}w},\label{eq:n2-JGm}\\
J(z)\,T(w)&\;\sim\;0.\label{eq:n2-JT}
\end{align}
\emph{(v) The supersymmetric OPE: $\Gp$--$\Gm$.}
\begin{equation}\label{eq:n2-GpGm}
\Gp(z)\,\Gm(w)\;\sim\;\frac{\ell}{(z{-}w)^3}
+\frac{J(w)}{(z{-}w)^2}
+\frac{T(w)+\frac12\partial J(w)}{z{-}w}.
\end{equation}
\emph{(vi) Vanishing OPEs.}
\begin{equation}\label{eq:n2-vanishing}
\Gp(z)\,\Gp(w)\;\sim\;0,\qquad
\Gm(z)\,\Gm(w)\;\sim\;0.
\end{equation}

\smallskip
\noindent\emph{Consistency check.}
There are $\binom{4}{2}+4=10$ independent pairs.  The non-trivial
zeroth products (simple-pole residues) are:
\begin{equation}\label{eq:n2-zeroth-products}
\begin{aligned}
\Gp_{(0)}\Gm &= T + \tfrac12\partial J, &\quad
J_{(0)}\Gp &= \Gp, &\quad
J_{(0)}\Gm &= -\Gm, \\
T_{(0)}\Gp &= \partial \Gp, &\quad
T_{(0)}\Gm &= \partial \Gm, &\quad
T_{(0)}J &= \partial J, \\
T_{(0)}T &= \partial T.
\end{aligned}
\end{equation}
All other zeroth products vanish:
$\Gp_{(0)}\Gp = \Gm_{(0)}\Gm = J_{(0)}J = J_{(0)}T = 0$.

The higher products relevant for the bar differential are:
\begin{equation}\label{eq:n2-higher-products}
\begin{aligned}
T_{(1)}T &= 2T, &\quad
T_{(1)}\Gp &= \tfrac32\,\Gp, &\quad
T_{(1)}\Gm &= \tfrac32\,\Gm, &\quad
T_{(1)}J &= J, \\
T_{(3)}T &= c/2, &\quad
\Gp_{(1)}\Gm &= J, &\quad
\Gp_{(2)}\Gm &= \ell, &\quad
J_{(1)}J &= \ell.
\end{aligned}
\end{equation}

\medskip
\noindent\textbf{Part III: The ordered bar complex at tensor
degree~$1$.}

${\Barch}^{\mathrm{ord}}_1
= \operatorname{span}\{s^{-1}T,\;s^{-1}\Gp,\;
s^{-1}\Gm,\;s^{-1}J\}$.
The four generators carry bar degrees:
\[
|s^{-1}T| = |s^{-1}J| = -1 \quad(\text{odd}),\qquad
|s^{-1}\Gp| = |s^{-1}\Gm| = 0 \quad(\text{even}).
\]
This grading controls all Koszul signs in the bar complex.
The conformal weights are $h(s^{-1}\Phi) = h(\Phi)$ (desuspension
does not shift conformal weight, only cohomological degree).

The $U(1)$-charge grading is:
$q(s^{-1}\Gp) = +1$, $q(s^{-1}\Gm) = -1$,
$q(s^{-1}T) = q(s^{-1}J) = 0$.

\medskip
\noindent\textbf{Part IV: Tensor degree~$2$, the $16$
ordered generators.}

${\Barch}^{\mathrm{ord}}_2
=\operatorname{span}\{s^{-1}\Phi_i\otimes s^{-1}\Phi_j
\mid 1\le i,j\le 4\}
\otimes C_\bullet(\mathrm{Conf}_2^{\mathrm{ord}}(\mathbb A^1))$.
Since $\mathrm{Conf}_2^{\mathrm{ord}}(\mathbb A^1)\simeq
\mathbb C^*$, only the $0$-chain (constant) and $1$-chain
(winding) sectors contribute.  We work in the constant sector
(which houses the bar differential).

The $16$ ordered generators, arranged by $U(1)$ charge
$q = q_i + q_j$:
\begin{center}
\renewcommand{\arraystretch}{1.15}
\begin{tabular}{cll}
\textbf{Charge $q$} & \textbf{Generators} & \textbf{Count} \\
\hline
$+2$ & $s^{-1}\Gp\otimes s^{-1}\Gp$ & $1$ \\
$+1$ & $s^{-1}\Gp\otimes s^{-1}T$,\;
$s^{-1}T\otimes s^{-1}\Gp$,\;
$s^{-1}\Gp\otimes s^{-1}J$,\;
$s^{-1}J\otimes s^{-1}\Gp$ & $4$ \\
$0$ & $s^{-1}T\otimes s^{-1}T$,\;
$s^{-1}J\otimes s^{-1}J$,\;
$s^{-1}T\otimes s^{-1}J$,\;
$s^{-1}J\otimes s^{-1}T$,\;
$s^{-1}\Gp\otimes s^{-1}\Gm$,\;
$s^{-1}\Gm\otimes s^{-1}\Gp$ & $6$ \\
$-1$ & $s^{-1}\Gm\otimes s^{-1}T$,\;
$s^{-1}T\otimes s^{-1}\Gm$,\;
$s^{-1}\Gm\otimes s^{-1}J$,\;
$s^{-1}J\otimes s^{-1}\Gm$ & $4$ \\
$-2$ & $s^{-1}\Gm\otimes s^{-1}\Gm$ & $1$
\end{tabular}
\end{center}

\emph{Koszul signs.}
For elements $\alpha = s^{-1}\Phi_i$ and
$\beta = s^{-1}\Phi_j$ in the tensor product, the
Koszul sign of the transposition
$\alpha\otimes\beta\mapsto\beta\otimes\alpha$ is
$(-1)^{|\alpha|\,|\beta|}$.  Since
$|s^{-1}T|=|s^{-1}J|=-1$ (odd) and
$|s^{-1}\Gp|=|s^{-1}\Gm|=0$ (even):
\begin{itemize}
\item bosonic--bosonic:
$s^{-1}T\otimes s^{-1}J \xmapsto{\text{flip}}
(-1)^{(-1)(-1)}s^{-1}J\otimes s^{-1}T
= -s^{-1}J\otimes s^{-1}T$
(anti-commuting in the unordered complex);
\item fermionic--fermionic:
$s^{-1}\Gp\otimes s^{-1}\Gm \xmapsto{\text{flip}}
(-1)^{(0)(0)}s^{-1}\Gm\otimes s^{-1}\Gp
= s^{-1}\Gm\otimes s^{-1}\Gp$
(commuting in the unordered complex);
\item mixed:
$s^{-1}J\otimes s^{-1}\Gp \xmapsto{\text{flip}}
(-1)^{(-1)(0)}s^{-1}\Gp\otimes s^{-1}J
= s^{-1}\Gp\otimes s^{-1}J$
(commuting in the unordered complex).
\end{itemize}
In the \emph{ordered} bar complex, all $16$ generators are
independent.  The Koszul signs above determine the
$R$-matrix at leading order
(Construction~\ref{constr:r-matrix-monodromy}).

\medskip
\noindent\textbf{Part V: The bar differential
$d_B\colon{\Barch}^{\mathrm{ord}}_2\to
{\Barch}^{\mathrm{ord}}_1$.}

The bar differential on a tensor-degree-$2$ element extracts
the \emph{collision residue}: for generators $\Phi_i$,
$\Phi_j$ of the chiral algebra,
\begin{equation}\label{eq:bar-diff-degree-2}
d_B(s^{-1}\Phi_i\otimes s^{-1}\Phi_j)
\;=\;
(-1)^{|s^{-1}\Phi_i|}\,
s^{-1}\bigl((\Phi_i)_{(0)}\Phi_j\bigr),
\end{equation}
where $(\Phi_i)_{(0)}\Phi_j$ is the zeroth product (residue of
the simple pole in the OPE).  The sign
$(-1)^{|s^{-1}\Phi_i|}$ is the standard Koszul sign from
the bar construction (the $d\log(z{-}w)$ kernel absorbs
one power of the pole).

Applying~\eqref{eq:bar-diff-degree-2}
using the zeroth products~\eqref{eq:n2-zeroth-products}:

\smallskip
\noindent\emph{Charge $q=0$ sector:}
\begin{align}
d_B(s^{-1}\Gp\otimes s^{-1}\Gm)
&\;=\;(-1)^0\,s^{-1}\bigl(\Gp_{(0)}\Gm\bigr)
\;=\;s^{-1}\bigl(T+\tfrac12\partial J\bigr),
\label{eq:dB-GpGm}\\
d_B(s^{-1}\Gm\otimes s^{-1}\Gp)
&\;=\;(-1)^0\,s^{-1}\bigl(\Gm_{(0)}\Gp\bigr).
\label{eq:dB-GmGp}
\end{align}
Now $\Gm_{(0)}\Gp = -\Gp_{(0)}\Gm + \partial(\Gm_{(1)}\Gp)$
by skew-symmetry of the $\lambda$-bracket.  Since
$\Gm_{(1)}\Gp = -\Gp_{(1)}\Gm = -J$
(from $\Gp_{(1)}\Gm = J$, with the sign from fermion
transposition in the skew-symmetry formula), we obtain
\begin{equation}\label{eq:Gm0Gp}
\Gm_{(0)}\Gp
\;=\;-(T+\tfrac12\partial J) + \partial(-J)
\;=\;-T-\tfrac32\partial J.
\end{equation}
Hence
$d_B(s^{-1}\Gm\otimes s^{-1}\Gp)
= s^{-1}(-T - \tfrac32\partial J)$.

\begin{align}
d_B(s^{-1}T\otimes s^{-1}T)
&\;=\;(-1)^{-1}\,s^{-1}(T_{(0)}T)
\;=\;-s^{-1}(\partial T),\label{eq:dB-TT}\\
d_B(s^{-1}T\otimes s^{-1}J)
&\;=\;(-1)^{-1}\,s^{-1}(T_{(0)}J)
\;=\;-s^{-1}(\partial J),\label{eq:dB-TJ}\\
d_B(s^{-1}J\otimes s^{-1}T)
&\;=\;(-1)^{-1}\,s^{-1}(J_{(0)}T)
\;=\;0,\label{eq:dB-JT}\\
d_B(s^{-1}J\otimes s^{-1}J)
&\;=\;(-1)^{-1}\,s^{-1}(J_{(0)}J)
\;=\;0.\label{eq:dB-JJ}
\end{align}

\smallskip
\noindent\emph{Charge $q=+1$ sector:}
\begin{align}
d_B(s^{-1}J\otimes s^{-1}\Gp)
&\;=\;(-1)^{-1}\,s^{-1}(J_{(0)}\Gp)
\;=\;-s^{-1}\Gp,\label{eq:dB-JGp}\\
d_B(s^{-1}T\otimes s^{-1}\Gp)
&\;=\;(-1)^{-1}\,s^{-1}(T_{(0)}\Gp)
\;=\;-s^{-1}(\partial\Gp),\label{eq:dB-TGp}\\
d_B(s^{-1}\Gp\otimes s^{-1}T)
&\;=\;(-1)^{0}\,s^{-1}(\Gp_{(0)}T).
\label{eq:dB-GpT}
\end{align}
Now $\Gp_{(0)}T$ is computed from skew-symmetry:
$\Gp_{(0)}T = T_{(0)}\Gp - \partial(T_{(1)}\Gp)
= \partial\Gp - \partial(\tfrac32\Gp) = -\tfrac12\partial\Gp$.
Skew-symmetry for mixed statistics requires care.
For $\Phi$ bosonic of bar degree $-1$ and $\Psi$ fermionic
of bar degree $0$, the OPE skew-symmetry reads
$\Psi_{(n)}T = -\sum_{j\ge 0}(-1)^{n+j}\tfrac{1}{j!}
\partial^j(T_{(n+j)}\Psi)$.
At $n=0$:
$\Gp_{(0)}T = -(T_{(0)}\Gp) + \partial(T_{(1)}\Gp)
= -\partial\Gp + \tfrac32\partial\Gp = \tfrac12\partial\Gp$.
Hence
$d_B(s^{-1}\Gp\otimes s^{-1}T)
= s^{-1}(\tfrac12\partial\Gp)$.
Similarly,
\begin{equation}\label{eq:dB-GpJ}
d_B(s^{-1}\Gp\otimes s^{-1}J)
\;=\;(-1)^{0}\,s^{-1}(\Gp_{(0)}J)
\;=\;-s^{-1}\Gp,
\end{equation}
where $\Gp_{(0)}J = -(J_{(0)}\Gp) + \partial(J_{(1)}\Gp)
= -\Gp + \partial(\ell\cdot 0) = -\Gp$
(since $J_{(1)}\Gp = 0$ because the $J\Gp$ OPE has only
a simple pole).

\smallskip
\noindent\emph{Charge $q=-1$ sector:}
By the $\Gp\leftrightarrow\Gm$, $q\mapsto -q$ symmetry:
\begin{align}
d_B(s^{-1}J\otimes s^{-1}\Gm)
&\;=\;s^{-1}\Gm,\label{eq:dB-JGm}\\
d_B(s^{-1}T\otimes s^{-1}\Gm)
&\;=\;-s^{-1}(\partial\Gm),\label{eq:dB-TGm}\\
d_B(s^{-1}\Gm\otimes s^{-1}T)
&\;=\;s^{-1}(\tfrac12\partial\Gm),\label{eq:dB-GmT}\\
d_B(s^{-1}\Gm\otimes s^{-1}J)
&\;=\;s^{-1}\Gm.\label{eq:dB-GmJ}
\end{align}

\smallskip
\noindent\emph{Charge $q=\pm 2$ sector:}
\begin{equation}\label{eq:dB-GpGp-GmGm}
d_B(s^{-1}\Gp\otimes s^{-1}\Gp)\;=\;0,\qquad
d_B(s^{-1}\Gm\otimes s^{-1}\Gm)\;=\;0,
\end{equation}
by the vanishing OPEs~\eqref{eq:n2-vanishing}.

\medskip
\noindent\textbf{Summary of $d_B$ on
${\Barch}^{\mathrm{ord}}_2$.}
The bar differential is a $16\times 4$ matrix.  In each
$U(1)$-charge sector:
\begin{center}
\renewcommand{\arraystretch}{1.3}
\begin{tabular}{ll}
$q=0$: &
$d_B = \begin{pmatrix}
\text{$\Gp\otimes\Gm$} \\
\text{$\Gm\otimes\Gp$} \\
\text{$T\otimes T$} \\
\text{$T\otimes J$} \\
\text{$J\otimes T$} \\
\text{$J\otimes J$}
\end{pmatrix}
\mapsto
\begin{pmatrix}
T+\frac12\partial J \\
-T-\frac32\partial J \\
-\partial T \\
-\partial J \\
0 \\
0
\end{pmatrix}$ \\[20pt]
$q=+1$: &
$d_B\colon\;
J\otimes\Gp\mapsto -\Gp,\;
T\otimes\Gp\mapsto -\partial\Gp,\;
\Gp\otimes T\mapsto \frac12\partial\Gp,\;
\Gp\otimes J\mapsto -\Gp$
\\[8pt]
$q=-1$: &
$d_B\colon\;
J\otimes\Gm\mapsto +\Gm,\;
T\otimes\Gm\mapsto -\partial\Gm,\;
\Gm\otimes T\mapsto \frac12\partial\Gm,\;
\Gm\otimes J\mapsto +\Gm$
\end{tabular}
\end{center}
(All outputs live in ${\Barch}^{\mathrm{ord}}_1$ with the
desuspension $s^{-1}$ suppressed.)

\medskip
\noindent\textbf{Part VI: The ordered coproduct
$\Delta\colon{\Barch}^{\mathrm{ord}}_2
\to {\Barch}^{\mathrm{ord}}_1\otimes
{\Barch}^{\mathrm{ord}}_1$.}

The coproduct is ordered deconcatenation, the same for
all chiral algebras, independent of the OPE:
\begin{equation}\label{eq:ordered-coproduct}
\Delta(s^{-1}\Phi_i\otimes s^{-1}\Phi_j)
\;=\;
(s^{-1}\Phi_i)\otimes(s^{-1}\Phi_j).
\end{equation}
This is the coassociative, non-cocommutative coproduct
encoding the $E_1$ (topological/associative) wing of the
$\SCchtop$ structure.  The reduced coproduct is
$\bar\Delta = \Delta$
on ${\Barch}^{\mathrm{ord}}_2$ (there is no lower-degree
contribution to subtract).

\emph{Compatibility with $d_B$.}
The Leibniz rule $\Delta\circ d_B
= (d_B\otimes\id + \id\otimes d_B)\circ\Delta$
holds on the nose: at tensor degree~$2$, both sides map to
${\Barch}^{\mathrm{ord}}_1$ (the left side via
$d_B$ followed by the trivial coproduct on degree~$1$,
which is $0$; the right side via $\Delta$ on degree~$2$
followed by $d_B$ on each tensor factor of degree~$1$,
which is also $0$ since $d_B=0$ on degree~$1$).  The
nontrivial Leibniz check begins at tensor degree~$3$.

\medskip
\noindent\textbf{Part VII: Verification that
$d_B^2 = 0$ (genus~$0$ uncurvedness).}

At genus $0$, the bar differential squares to zero.  We
verify this at tensor degree~$3$: the bar differential
$d_B\colon{\Barch}^{\mathrm{ord}}_3\to
{\Barch}^{\mathrm{ord}}_2$ has two face maps
$d_B = d_1 - d_2$, where
\[
d_1(s^{-1}a\otimes s^{-1}b\otimes s^{-1}c)
\;=\;
(-1)^{|s^{-1}a|}\,s^{-1}(a_{(0)}b)\otimes s^{-1}c,
\]
\[
d_2(s^{-1}a\otimes s^{-1}b\otimes s^{-1}c)
\;=\;
(-1)^{|s^{-1}a|+|s^{-1}b|}\,
s^{-1}a\otimes s^{-1}(b_{(0)}c).
\]
Then $d_B^2 = d_B\circ d_B$ on degree~$3$ elements:
the composition $d_B^{(1)}\circ d_B^{(2)}$ (applied
to a degree-$3$ element, first reducing to degree~$2$,
then to degree~$1$) gives
\begin{align*}
d_B^2(s^{-1}a\otimes s^{-1}b\otimes s^{-1}c)
&\;=\;
\pm s^{-1}\bigl((a_{(0)}b)_{(0)}c\bigr)
\mp s^{-1}\bigl(a_{(0)}(b_{(0)}c)\bigr) \\
&\;=\;
\pm s^{-1}\bigl([a_{(0)},b_{(0)}]c
- a_{(0)}(b_{(0)}c) + (a_{(0)}b)_{(0)}c\bigr).
\end{align*}
The Borcherds identity (Jacobi for vertex algebras) at
$m=n=0$ gives exactly
$(a_{(0)}b)_{(0)}c = a_{(0)}(b_{(0)}c) - b_{(0)}(a_{(0)}c)$
for generators where the sign conventions are tracked
properly.  Therefore $d_B^2 = 0$.  This is automatic and
holds for \emph{any} vertex algebra at genus~$0$; the
verification merely confirms that our sign conventions are
consistent.

\emph{The genus-$g$ curvature.}
At genus $g\ge 1$, the curved bar complex has
$d_B^2 = \kappa(A)\cdot\omega_g$ where
$\kappa(A)$ is the curvature of $A$ and $\omega_g$ is the
Hodge class.  For the $\mathcal N=2$ SCA, the curvature is
$\kappa = (k+4)/4$ (see Vol~I Master Table; note this is \emph{not} $c/24$).

\medskip
\noindent\textbf{Part VIII: The ghost central charge from the
bar complex trace.}

The DS reduction
$\mathrm{DS}_{f_{\mathrm{sub}}}\colon
\wsl_3{}_k \to W_k(\mathfrak{sl}_3,f_{\mathrm{sub}})$
removes the roots in the nilradical of the
Jacobson--Morozov parabolic $\mathfrak p_{JM}
\subset\mathfrak{sl}_3$.  For the subregular nilpotent
$f_{\mathrm{sub}}$ in $\mathfrak{sl}_3$, the
Jacobson--Morozov $\mathfrak{sl}_2$-triple
$(e_{\mathrm{sub}}, h_{\mathrm{sub}}, f_{\mathrm{sub}})$
determines the grading of $\mathfrak{sl}_3$ by
$\mathrm{ad}(h_{\mathrm{sub}})$-eigenvalues.

The root decomposition of $\mathfrak{sl}_3$ under
$\mathrm{ad}(h_{\mathrm{sub}})$ has:
\begin{itemize}
\item grade $0$: a $3$-dimensional subalgebra
$\mathfrak{sl}_2\oplus\mathbb C$ (the Levi factor of
$\mathfrak p_{JM}$),
\item grade $+1$: a $2$-dimensional space (two positive roots),
\item grade $-1$: a $2$-dimensional space (two negative roots).
\end{itemize}
The BRST complex introduces $bc$-ghosts for the grade-$\ge 1$
directions: two ghost pairs $(b_\alpha, c_\alpha)$,
$\alpha=1,2$.  The conformal weights of these ghosts
\emph{depend on the level $k$} because the JM grading
interacts with the Sugawara construction.  Specifically,
the ghost conformal weights are determined by the
$\mathrm{ad}(h_{\mathrm{sub}})$-eigenvalues $n_\alpha$
and the conformal weights $h_\alpha$ of the
corresponding root vectors:
$h(c_\alpha) = 1 - h_\alpha$, $h(b_\alpha) = h_\alpha$.
Since $h_\alpha$ depends on the level (through the
Sugawara denominator $k+h^\vee = k+3$), so does the
ghost central charge.

The ghost central charge is computed from the DS formula:
\begin{equation}\label{eq:c-ghost-subregular}
c_{\mathrm{ghost}}
\;=\;
c(\wsl_3{}_k) - c(W_k(\mathfrak{sl}_3,f_{\mathrm{sub}}))
\;=\;
\frac{8k}{k+3} - \frac{3(2k+1)}{k+3}
\;=\;
\frac{2k-3}{k+3}.
\end{equation}

\emph{Level-dependence.}
Unlike the principal case where
$c_{\mathrm{ghost}} = 2(N^2-1) - c(W_k^{\mathrm{prin}})
= N(N-1)$ (a constant independent of $k$), the subregular
ghost central charge~\eqref{eq:c-ghost-subregular} depends
on $k$.  This is the central new phenomenon.

\emph{Bar complex interpretation.}
The ghost central charge appears as the trace of the bar
complex:
\begin{equation}\label{eq:c-ghost-bar-trace}
c_{\mathrm{ghost}}
\;=\;
\operatorname{str}_{{\Barch}^{\mathrm{ord}}_1}
\bigl(L_0\bigr)
-
\operatorname{str}_{{\Barch}^{\mathrm{ord}}_1(\wsl_3)}
\bigl(L_0\bigr),
\end{equation}
where the supertrace is over the degree-$1$ bar generators,
weighted by conformal weight.
For $\wsl_3{}_k$: generators $e_{ij}$ ($i\neq j$) and
$h_i$ ($i=1,2$), all of weight $1$, all bosonic, giving
$\operatorname{str}(L_0) = 8\cdot 1 = 8$ (normalised by
$1/(k+3)$ from the Sugawara construction, giving
$8/(k+3)$ per unit of central charge).
For $W_k(\mathfrak{sl}_3,f_{\mathrm{sub}})$: generators
$J$ (weight $1$, bosonic, $\operatorname{str}=+1$),
$\Gp$, $\Gm$ (weight $3/2$, fermionic,
$\operatorname{str}=-3/2$ each),
$T$ (weight $2$, bosonic, $\operatorname{str}=+2$).
Total: $1 - 3/2 - 3/2 + 2 = 1$.
The difference in the normalised supertrace gives
$c_{\mathrm{ghost}} = c_{\mathrm{aff}} - c_W
= 8k/(k+3) - 3(2k+1)/(k+3) = (2k-3)/(k+3)$.

\emph{Special values.}
$c_{\mathrm{ghost}}=0$ at $k=3/2$ (the ghost system decouples);
$c_{\mathrm{ghost}}=4$ at $k=-15/2$
($=-h^\vee - 3/2 = -3 - 3/2$, the ``additional critical
level'');
$c_{\mathrm{ghost}}\to 2$ as $k\to\infty$ (the
semi-classical limit, matching $\dim(\mathfrak n_+^{JM})=2$,
consistent with the naive $bc$-ghost count in the classical
limit).

\medskip
\noindent\textbf{Part IX: The Koszul dual, subregular
to minimal.}

\emph{Statement.}
The Koszul dual of
$A = W_k(\mathfrak{sl}_3,f_{\mathrm{sub}})$ is conjectured
to be
\begin{equation}\label{eq:koszul-dual-conjecture}
A^!\;\stackrel{?}{\simeq}\;
W_{k^\vee}(\mathfrak{sl}_3,f_{\mathrm{sub}}^t),
\qquad
k^\vee = -k - 2h^\vee = -k-6,
\end{equation}
where $f_{\mathrm{sub}}^t$ is the \emph{transpose
nilpotent}.  In $\mathfrak{sl}_3$, the transpose of the
subregular nilpotent is the minimal nilpotent:
$f_{\mathrm{sub}}^t = f_{\min}$.  Hence:
\begin{equation}\label{eq:koszul-dual-explicit}
W_k(\mathfrak{sl}_3,f_{\mathrm{sub}})^!
\;\stackrel{?}{\simeq}\;
W_{-k-6}(\mathfrak{sl}_3,f_{\min}).
\end{equation}
This is a Koszul duality between
\emph{different $W$-algebras}: the subregular
$\mathcal N=2$ SCA (with generators of weights
$1,\frac32,\frac32,2$) and the minimal $W$-algebra
$W_{-k-6}(\mathfrak{sl}_3,f_{\min})$, which is
the Bershadsky--Polyakov algebra (with generators of weights
$1,\frac32,\frac32,2$, the same weight spectrum but
different OPE coefficients).

\emph{Evidence from central charges.}
The Feigin--Frenkel duality at the affine level gives
$c(\wsl_3{}_k) + c(\wsl_3{}_{-k-6}) = 16$
(since $c(\wsl_3{}_k) = 8k/(k+3)$ and
$c(\wsl_3{}_{-k-6}) = 8(-k-6)/(-k-6+3) = -8(k+6)/(k+3)$,
sum $= (8k-8k-48)/(k+3)=-48/(k+3)$;
corrected: $c(\wsl_3{}_{k'}) = 8k'/(k'+3)$ with
$k'=-k-6$, so $k'+3=-k-3$, giving
$c = 8(-k-6)/(-k-3) = 8(k+6)/(k+3)$,
and $c(\wsl_3{}_k)+c(\wsl_3{}_{k'})
= 8k/(k+3)+8(k+6)/(k+3) = 8(2k+6)/(k+3) = 16(k+3)/(k+3)
= 16$).

The W-algebra central charges:
$c(A) = 3(2k+1)/(k+3)$ and
$c(A^!) = c(W_{-k-6}(\mathfrak{sl}_3,f_{\min}))$.
The Bershadsky--Polyakov central charge at level $k'$ is
$c_{BP}(k') = (k'-1)(3k'-4)/(k'+3/2) - 1/2$; at
$k'=-k-6$ this requires a careful substitution.

The correct approach uses the complementarity identity.
For affine KM algebras, $\kappa + \kappa^! = 0$.
The subregular reduction modifies this.  The relevant
check is whether the complementarity sum
\begin{equation}\label{eq:complementarity-check}
c(A) + c(A^!) \;=\; c_{\mathrm{self\text{-}dual}}
\end{equation}
equals the self-dual central charge.

\emph{Evidence from pole structures.}
Both $A$ and $A^!$ have four generators of weights
$1,\frac32,\frac32,2$.  The maximal pole orders in the
OPEs are:
\begin{center}
\renewcommand{\arraystretch}{1.1}
\begin{tabular}{lcc}
\textbf{OPE pair} & \textbf{$A$ (SCA)} & \textbf{$A^!$ (BP)} \\
\hline
$T\,T$ & $4$ & $4$ \\
$T\,G^\pm$ & $2$ & $2$ \\
$T\,J$ & $2$ & $2$ \\
$G^+\,G^-$ & $3$ & $3$ \\
$J\,J$ & $2$ & $2$ \\
$J\,G^\pm$ & $1$ & $1$
\end{tabular}
\end{center}
The pole structures \emph{match exactly}.  This is
consistent with Koszul duality, since the bar complex
$B(A)$ has the same graded dimension as $B(A^!)$ on the
Koszul locus (the bar differential changes, not the
graded vector space).

\emph{Evidence from the bar differential.}
The bar kernel absorbs one pole: the collision
residue $r_n$ comes from the $(n+1)$-st order pole.
Therefore $d_B$ at tensor degree~$2$ encodes the zeroth
products $\Phi_{(0)}\Psi$ (from the simple pole),
which for both $A$ and $A^!$ are:
\begin{itemize}
\item $G^+_{(0)}G^-$: a weight-$2$ bosonic composite
($T + \frac12\partial J$ for the SCA,
$T + \alpha\partial J$ for BP with $\alpha$ depending
on conventions);
\item $J_{(0)}G^\pm = \pm G^\pm$ (the $U(1)$ charge is the
same for both algebras);
\item $T_{(0)}\Phi = \partial\Phi$ (the Virasoro translation).
\end{itemize}
The \emph{differences} between the two bar differentials
arise from the OPE coefficients at higher poles: the
second-order pole coefficient $\Gp_{(1)}\Gm = J$ for the
SCA but $\Gp_{(1)}\Gm = J'$ (a shifted current) for BP.
These differences manifest in the $R$-matrix at order
$1/z$ and beyond.

\emph{Obstruction to proving the conjecture.}
The conjecture~\eqref{eq:koszul-dual-explicit} requires:
\begin{enumerate}[label=(\alph*)]
\item Koszulness of $A=W_k(\mathfrak{sl}_3,f_{\mathrm{sub}})$
: this follows from the one-loop exactness criterion
  (Theorem~\ref{thm:one-loop-koszul}) if the subregular
  DS reduction preserves one-loop exactness;
\item the DS/KD intertwining
  (Theorem~\ref{conj:DS}) for the non-principal case;
\item an explicit identification of the Verdier dual
  $D_{\mathrm{Ran}}(B(A))$ with $B(A^!)$.
\end{enumerate}
Obstruction~(b) is exactly the non-principal corridor
frontier
(Conjecture~\ref*{V1-conj:ds-kd-arbitrary-nilpotent}).
This computation provides the first explicit data point for
that conjecture.

\medskip
\noindent\textbf{Part X: Higher shadows and the shadow tower.}

The $\mathcal N=2$ SCA has shadow depth~$\infty$:
the fermionic generators $G^\pm$ produce nonvanishing
transferred operations at all arities via the $G^+G^-$ OPE
(which has a simple pole producing $T+\frac12\partial J$).
The quartic contact shadow involves the composite field
${:}\Gp\Gm{:}$ and the cubic shadow involves $T\cdot J$.
The shadow tower for the $\mathcal N=2$ SCA is strictly
richer than for the Virasoro algebra: it has a
$4$-dimensional deformation space
($\partial_c T$, $\partial_c \Gp$, $\partial_c \Gm$,
$\partial_c J$) and the genus-$1$ Hessian is a
$4\times 4$ matrix.

\medskip
\noindent\textbf{Part XI: The $R$-matrix and spectral structure.}

The $R$-matrix
(Construction~\ref{constr:r-matrix-monodromy}) for the
$\mathcal N=2$ SCA acts on
$V\otimes V$ where $V = \operatorname{span}\{T,\Gp,\Gm,J\}$.
At leading order:
\begin{equation}\label{eq:r-matrix-n2-sca}
R(z) \;=\; \id + r_0/z + r_1/z^2 + r_2/z^3 + \cdots,
\end{equation}
where $r_0$ is the classical $r$-matrix built from zeroth
products:
\begin{equation}\label{eq:classical-r-n2}
r_0 \;=\; \sum_{i,j} (\Phi_i)_{(0)}(\Phi_j)\;\cdot\;
e_i\otimes e_j.
\end{equation}
The non-zero entries of $r_0$ in the basis
$(T,\Gp,\Gm,J)$ are:
\[
r_0 \;=\;
\begin{pmatrix}
\partial T & \frac12\partial\Gp & \frac12\partial\Gm & \partial J \\
-\partial\Gp & 0 & T+\frac12\partial J & -\Gp \\
-\partial\Gm & -T-\frac32\partial J & 0 & \Gm \\
0 & -\Gp & \Gm & 0
\end{pmatrix},
\]
where each entry $(r_0)_{ij} = (\Phi_i)_{(0)}\Phi_j$.
The higher-order terms $r_1$, $r_2$ involve the first and
second products $\Phi_{(1)}\Psi$, $\Phi_{(2)}\Psi$:
\begin{equation}\label{eq:r1-n2}
r_1 \;=\;
\begin{pmatrix}
2T & \frac32\Gp & \frac32\Gm & J \\
0 & 0 & J & 0 \\
0 & -J & 0 & 0 \\
0 & 0 & 0 & \ell
\end{pmatrix},
\qquad
r_2 \;=\;
\begin{pmatrix}
0 & 0 & 0 & 0 \\
0 & 0 & \ell & 0 \\
0 & 0 & 0 & 0 \\
0 & 0 & 0 & 0
\end{pmatrix}.
\end{equation}
The entry $(r_2)_{\Gp,\Gm} = \ell = c/3$ reflects the
third-order pole in the $\Gp\Gm$ OPE (after bar-kernel
absorption: a third-order pole in the OPE contributes
at $1/z^3$ in the $R$-matrix, i.e.\ at $r_2$).

The $R$-matrix satisfies the Yang--Baxter equation
$R_{12}(z_1{-}z_2)\,R_{13}(z_1{-}z_3)\,R_{23}(z_2{-}z_3)
= R_{23}(z_2{-}z_3)\,R_{13}(z_1{-}z_3)\,R_{12}(z_1{-}z_2)$
by $d_B^2=0$ on ${\Barch}^{\mathrm{ord}}_3$
(Proposition~\ref{prop:ybe-from-d-squared}).

\emph{The $U(1)$ decomposition.}
The $R$-matrix decomposes into $U(1)$-charge sectors:
$R(z) = \bigoplus_{q} R^{(q)}(z)$, where $R^{(q)}$ acts on
the charge-$q$ subspace of $V\otimes V$.  This decomposition
is exact (not merely approximate) because $[J_0, R(z)] = 0$:
the $R$-matrix preserves the total $U(1)$ charge.  This
$U(1)$ decomposition has no analogue for the Virasoro algebra
or for principal $W$-algebras without a $U(1)$ current.

\emph{Comparison with the affine case.}
For $\wsl_2{}_k$, the classical $r$-matrix is the split
Casimir $r_0 = \Omega/z$ with $\Omega\in\mathfrak{sl}_2
\otimes\mathfrak{sl}_2$.  For the $\mathcal N=2$ SCA, the
classical $r$-matrix~\eqref{eq:classical-r-n2} is not a
Casimir element of any finite-dimensional Lie algebra:
the presence of $\partial$-derivatives and the mixing of
different conformal weights prevent a Casimir
interpretation.  This is the hallmark of a $W$-algebra
(as opposed to a Kac--Moody algebra) at the level of
the ordered bar complex.
\end{computation}

\subsection{The Bershadsky--Polyakov algebra:
the ordered bar complex of a minimal $W$-algebra}
\label{subsec:ordered-bar-bp}

\providecommand{\Gp}{G^+}
\providecommand{\Gm}{G^-}
\providecommand{\wsl}{\widehat{\mathfrak{sl}}}

The preceding computation treated the subregular
$W_k(\mathfrak{sl}_3,f_{\mathrm{sub}})$, which is the
$\mathcal N=2$ SCA\@.  Its conjectural Koszul dual
\eqref{eq:koszul-dual-explicit} is the minimal $W$-algebra
$W_{k'}(\mathfrak{sl}_3,f_{\min})$, the
\emph{Bershadsky--Polyakov algebra} $\mathcal{B}^{k'}$.
We now compute the ordered bar complex of $\mathcal{B}^{k'}$
directly.  This is the first ordered bar computation for a
\emph{minimal} (as opposed to subregular or principal)
$W$-algebra, and it exhibits four phenomena absent from all
previous examples:
\begin{enumerate}[label=(\roman*)]
\item a four-level depth stratification
(gravity/matter/current/charge) reflecting the
$\mathfrak{sl}_3$ Casimir structure;
\item a charge asymmetry in the collision residue
($r^{\mathrm{coll}}_{\Gp\Gm}\neq -r^{\mathrm{coll}}_{\Gm\Gp}$
up to sign) that distinguishes creation from annihilation;
\item non-primitive coproduct on the stress tensor, with
cross-term $J\otimes J$ from the surviving weight-$1$
generator (cf.\ Remark~\ref{rem:bp-coproduct-resolved});
\item selective shadow propagation: the quartic (gravitational)
shadow in the $(T,T)$ sector does \emph{not} propagate into
the matter sector at $m_3$ level.
\end{enumerate}

\begin{computation}[The ordered bar complex of the
Bershadsky--Polyakov algebra
$\mathcal{B}^{k'} = W_{k'}(\mathfrak{sl}_3,f_{\min})$;
\ClaimStatusProvedHere]
\label{comp:ordered-bar-bp}
\index{Bershadsky--Polyakov algebra!ordered bar complex|textbf}
\index{bar complex!ordered!Bershadsky--Polyakov|textbf}
\index{W-algebra@$\mathcal W$-algebra!minimal!ordered bar|textbf}

\medskip
\noindent\textbf{Part~I: Generators and OPE data.}

The algebra $\mathcal{B}^{k'}$ has four strong generators:
\begin{center}
\renewcommand{\arraystretch}{1.15}
\begin{tabular}{lccc}
\textbf{Generator} & \textbf{Conf.\ weight} &
\textbf{Statistics} & \textbf{$U(1)$ charge} \\
\hline
$T$ & $2$ & bosonic & $0$ \\
$\Gp$ & $3/2$ & fermionic & $+1$ \\
$\Gm$ & $3/2$ & fermionic & $-1$ \\
$J$ & $1$ & bosonic & $0$
\end{tabular}
\end{center}
The weight spectrum coincides with that of the $\mathcal N=2$
SCA; the two algebras are distinguished by their OPE
coefficients.

\emph{Central charge.}
The Bershadsky--Polyakov central charge at level $k'$ in
the manuscript's convention is
\begin{equation}\label{eq:bp-central-charge}
c_{\mathrm{BP}}(k')
\;=\;
\frac{k'-15}{k'+3}.
\end{equation}
This differs from the Adamovi\'{c}--Arakawa convention
$c_{\mathrm{AA}}(k')=-6(k'+2)(k'-1)/(k'+3)$;
the two are related by a shift in the meaning of the level
parameter.  In the manuscript's convention, the
complementarity identity reads
$c_{\mathrm{BP}}(k')+c_{\mathrm{BP}}(-k'-6) = 2$
(cf.\ \eqref{eq:complementarity-check}),
confirmed by direct substitution:
$(k'-15)/(k'+3)+(-k'-21)/(-k'-3)
=(k'-15)/(k'+3)+(k'+21)/(k'+3)
=36/(k'+3)\neq 2$ --- corrected:
the identity uses the \emph{SCA} central charge
$c_{\mathrm{SCA}}(k)=3(2k+1)/(k+3)$ on one side and the
BP charge on the other.  The correct duality pair is
\begin{equation}\label{eq:bp-complementarity}
c_{\mathrm{SCA}}(k) + c_{\mathrm{BP}}(-k-6)
\;=\; 2
\end{equation}
for $k'=-k-6$, matching
Proposition~\texttt{prop:partition-dependent-complementarity}
of Volume~I\@.

\medskip
\noindent\textbf{Part~II: The complete $m_2$ table.}

The $4\times 4$ generator space produces $16$ ordered pairs.
Five follow from skew-symmetry (with Koszul sign):
\begin{equation}\label{eq:bp-skew-symmetry}
\{b_\lambda\, a\}
\;=\;
-(-1)^{|a||b|}\,\{a_{-\lambda-\partial}\,b\},
\end{equation}
where $|a|$ denotes the parity ($0$ for bosonic, $1$ for
fermionic).  The remaining $11$ independent $\lambda$-brackets
(from which the collision residues and bar differential are
extracted) are:

\begin{align}
\{T_\lambda\, T\}
&\;=\;
\tfrac{c_{\mathrm{BP}}}{2}\,\lambda^3
+ 2T\lambda + \partial T,
\label{eq:bp-TT}
\\[4pt]
\{T_\lambda\, \Gp\}
&\;=\;
\tfrac{3}{2}\,\Gp\lambda + \partial\Gp,
\label{eq:bp-TGp}
\\[4pt]
\{T_\lambda\, \Gm\}
&\;=\;
\tfrac{3}{2}\,\Gm\lambda + \partial\Gm,
\label{eq:bp-TGm}
\\[4pt]
\{T_\lambda\, J\}
&\;=\;
J\lambda + \partial J,
\label{eq:bp-TJ}
\\[4pt]
\{J_\lambda\, J\}
&\;=\;
K\lambda,
\label{eq:bp-JJ}
\\[4pt]
\{J_\lambda\, \Gp\}
&\;=\;
\Gp,
\label{eq:bp-JGp}
\\[4pt]
\{J_\lambda\, \Gm\}
&\;=\;
-\Gm,
\label{eq:bp-JGm}
\\[4pt]
\{\Gp_\lambda\, \Gm\}
&\;=\;
K\lambda^2 + J\lambda
+ T + \tfrac{3}{2K}{:}JJ{:} + \tfrac{1}{2}\partial J,
\label{eq:bp-GpGm}
\\[4pt]
\{\Gm_\lambda\, \Gp\}
&\;=\;
K\lambda^2 - J\lambda
+ T + \tfrac{3}{2K}{:}JJ{:} - \tfrac{1}{2}\partial J,
\label{eq:bp-GmGp}
\\[4pt]
\{\Gp_\lambda\, \Gp\}
&\;=\; 0,
\label{eq:bp-GpGp}
\\[4pt]
\{\Gm_\lambda\, \Gm\}
&\;=\; 0.
\label{eq:bp-GmGm}
\end{align}
Here $K=k'+3$ is the shifted level
(so that the $J$--$J$ OPE pole is $K/(z-w)^2$),
and ${:}JJ{:}$ denotes the normal-ordered product.
The critical entry is
\eqref{eq:bp-GpGm}:
the $\Gp$--$\Gm$ bracket contains the
\emph{Sugawara-like composite} $T+\frac{3}{2K}{:}JJ{:}$, which
is the modified stress tensor of the coset
$\wsl_3{}_k / \wsl_2{}_{k'}$.

\emph{Charge asymmetry.}
The collision residue (applying AP19: the $d\log$ kernel absorbs
one pole) extracts
\begin{equation}\label{eq:bp-collision-residue-asymmetry}
r^{\mathrm{coll}}_{\Gp\Gm}(z)
\;=\;
\frac{K}{z^2} + \frac{J}{z},
\qquad
r^{\mathrm{coll}}_{\Gm\Gp}(z)
\;=\;
\frac{K}{z^2} - \frac{J}{z}.
\end{equation}
The quadratic pole coefficient $K$ is symmetric, but the
simple pole coefficient $\pm J$ is \emph{antisymmetric}:
it is the $U(1)$ charge operator itself.  This asymmetry
has no analogue in the Virasoro or principal $\mathcal{W}_N$
algebras (which have no weight-$1$ generator) and is the
collision-residue avatar of charge-conjugation asymmetry.
For the SCA (Computation~\ref{comp:ordered-bar-n2-sca}),
the analogous entries are
$r^{\mathrm{coll}}_{\Gp\Gm}=K'/z^2+J'/z$ and
$r^{\mathrm{coll}}_{\Gm\Gp}=K'/z^2-J'/z$ with
different constants $K'$, $J'$; the \emph{structural}
asymmetry is identical.

\medskip
\noindent\textbf{Part~III: Four-level depth stratification.}

\begin{remark}[Depth stratification of the BP ordered bar
complex]
\label{rem:bp-depth-stratification}
\index{Bershadsky--Polyakov algebra!depth stratification|textbf}
\index{depth!stratification!Bershadsky--Polyakov}
The maximal pole order in each OPE sector, after bar-kernel
absorption (AP19), determines the depth of the corresponding
component of the ordered bar differential:
\begin{center}
\renewcommand{\arraystretch}{1.15}
\begin{tabular}{llcc}
\textbf{Sector} & \textbf{Generators}
& \textbf{Max OPE pole} & \textbf{Depth} \\
\hline
Gravity & $(T,T)$ & quartic & $3$ \\
Matter & $(\Gp,\Gm)$ & cubic & $2$ \\
Current & $(J,J)$ & double & $1$ \\
Charge & $(J,G^\pm)$ & simple & $0$
\end{tabular}
\end{center}
The depth column records the maximal power of $1/z$ in the
collision residue (pole order minus one, by AP19).  The
stratification is:
\begin{itemize}
\item \textbf{Gravity sector} (depth $3$): The $T$--$T$ OPE
is a quartic pole, identical in structure to the Virasoro
algebra.  This sector alone is class~$\mathbf{M}$
(infinite shadow tower).
\item \textbf{Matter sector} (depth $2$): The $\Gp$--$\Gm$
OPE has a cubic pole (the $K\lambda^2$ term).  The collision
residue has depth $2$, placing this sector in
class~$\mathbf{C}$ (finite but nontrivial shadow tower).
\item \textbf{Current sector} (depth $1$): The $J$--$J$ OPE
has a double pole ($K\lambda$).  The collision residue is
$K/z$, placing this sector in class~$\mathbf{L}$
(formal, Heisenberg-type).
\item \textbf{Charge sector} (depth $0$): The $J$--$G^\pm$
OPEs are simple poles ($\pm G^\pm$, no $\lambda$-dependence).
The collision residue is a constant, placing this sector in
class~$\mathbf{G}$ (abelian).
\end{itemize}
The BP algebra is therefore \emph{stratified across all four
depth classes}: $\mathbf{G}$, $\mathbf{L}$, $\mathbf{C}$,
$\mathbf{M}$.  No previous example in the manuscript exhibits
all four classes simultaneously.  The principal $W$-algebras
are pure class~$\mathbf{M}$ (all pole orders quartic
for Virasoro, or mixed quartic/lower for $\mathcal{W}_N$ but
with no class~$\mathbf{G}$ sector).  The Heisenberg algebra
is pure class~$\mathbf{L}$.  The Kac--Moody algebras are
class~$\mathbf{L}$ (double pole only).  The BP algebra is
the first example exhibiting the complete four-level hierarchy.
\end{remark}

\medskip
\noindent\textbf{Part~IV: The $m_3$ operations.}

The transferred ternary products
$m_3\colon (s^{-1}\overline{\mathcal{B}})^{\otimes 3}
\to s^{-1}\overline{\mathcal{B}}$ are computed by
evaluating the bar differential $d_B$ on tensor degree~$3$
elements and projecting.  We record the key components.

\emph{(a) $m_3(T,T,T)$: the Virasoro sub-PVA.}
The gravity sector closes on itself at all arities.  The
restriction of the $\Ainf$-structure to the
$1$-dimensional subspace $\langle T\rangle$ yields the
Virasoro $\Ainf$-algebra at $c=c_{\mathrm{BP}}$.
In particular,
\begin{equation}\label{eq:bp-m3-TTT}
m_3(T,T,T)
\;=\;
m_3^{\mathrm{Vir}}(T,T,T)\big|_{c=c_{\mathrm{BP}}},
\end{equation}
the same expression as for the Virasoro algebra
(Computation~\texttt{comp:virasoro-m3} of the examples chapter),
with $c$ specialised to $c_{\mathrm{BP}}(k')=(k'-15)/(k'+3)$.
The $T$-sector shadow tower is therefore the Virasoro shadow
tower at the BP central charge: infinite depth,
class~$\mathbf{M}$.

\emph{(b) $m_3(T,T,\Gp)$: quartic shadow does not propagate
into the matter sector.}
The quartic pole $c_{\mathrm{BP}}\lambda^3/2$ in the
$T$--$T$ OPE generates a depth-$3$ collision residue.
However, when one of the three inputs is a fermionic
generator $\Gp$ (or $\Gm$), the transferred product
$m_3(T,T,\Gp)$ receives contributions only from the
$T$--$\Gp$ collision (which has a double pole,
depth~$1$) and from the composite
$m_2(m_2(T,T),\Gp) - m_2(T,m_2(T,\Gp))$ (the
$\Ainf$-relation).  The quartic pole in the $T$--$T$
channel produces an intermediate element proportional to
$T$ (after $d_B$ acts), and the subsequent
$T$--$\Gp$ pairing sees only the quadratic pole.  Thus:
\begin{equation}\label{eq:bp-m3-TTGp}
m_3(T,T,\Gp)
\;=\;
\bigl(\text{depth-$1$ contribution}\bigr)
\;\in\;
s^{-1}\langle\Gp,\partial\Gp,\partial^2\Gp\rangle.
\end{equation}
The quartic shadow is \emph{confined to the gravity sector}:
it does not propagate into the matter sector at $m_3$ level.
This is a selection rule from the depth stratification ---
the $T$--$G^\pm$ OPE has insufficient pole order to transmit
the full gravitational shadow.

\emph{(c) $m_3(G^\pm,G^\pm,X)=0$: fermion nilpotency.}
For any generator $X\in\{T,\Gp,\Gm,J\}$:
\begin{equation}\label{eq:bp-m3-fermion-nil}
m_3(\Gp,\Gp,X) \;=\; 0,
\qquad
m_3(\Gm,\Gm,X) \;=\; 0.
\end{equation}
This vanishing is forced by the Koszul sign rule on the
desuspended bar complex: $s^{-1}\Gp$ is \emph{bosonic}
(the desuspension shifts parity), and the skew-symmetry
$\{\Gp_\lambda\Gp\}=0$ of the OPE
\eqref{eq:bp-GpGp} ensures that the two $\Gp$-inputs
contribute zero at every stage of the HPL iteration.
The result holds at all arities:
$m_n(\Gp,\Gp,X_1,\dots,X_{n-2})=0$ for all $n\ge 2$.

\emph{(d) $m_3(J,J,J)=0$: Heisenberg formality.}
The current sector is class~$\mathbf{L}$ (formal):
\begin{equation}\label{eq:bp-m3-JJJ}
m_3(J,J,J) \;=\; 0.
\end{equation}
The $J$--$J$ OPE has a double pole only.  The bar
differential at tensor degree~$3$ on
$(s^{-1}J)^{\otimes 3}$ receives a contribution from
the double pole in the $(J,J)$ collision, producing
$K\cdot 1$ (a scalar), which then pairs with $J$ via
$m_2(K\cdot 1,J)=0$ (the unit acts trivially on the
reduced bar complex).  All higher $m_n(J,\dots,J)$
vanish by the same mechanism: this is exactly the formality
of the Heisenberg algebra (class~$\mathbf{L}$) embedded
inside the BP algebra.

\medskip
\noindent\textbf{Part~V: Non-primitive coproduct.}

\begin{proposition}[BP coproduct structure]
\label{prop:bp-coproduct}
\index{Bershadsky--Polyakov algebra!coproduct|textbf}
\index{coproduct!non-primitive!Bershadsky--Polyakov}
The linear transferred coproduct
$\Delta_{z,1}\colon s^{-1}\overline{\mathcal{B}}
\to (s^{-1}\overline{\mathcal{B}})^{\otimes 2}$
on the ordered bar complex of $\mathcal{B}^{k'}$ satisfies:
\begin{enumerate}[label=\textup{(\alph*)}]
\item $\Delta_{z,1}(\Gp)
= \Gp\otimes 1 + 1\otimes\Gp$
\quad\textup{(primitive)};
\item $\Delta_{z,1}(\Gm)
= \Gm\otimes 1 + 1\otimes\Gm$
\quad\textup{(primitive)};
\item $\Delta_{z,1}(J)
= J\otimes 1 + 1\otimes J$
\quad\textup{(primitive)};
\item $\Delta_{z,1}(T)
= T\otimes 1 + 1\otimes T
+ \frac{1}{z}\,J\otimes J$
\quad\textup{(non-primitive)}.
\end{enumerate}
\end{proposition}

\begin{proof}
Items (a)--(c): the generators $\Gp$, $\Gm$, $J$ are
\emph{fundamental} --- they arise directly from the
$\mathfrak{sl}_3$ root/Cartan currents, not from
composites.  The split Casimir
$\Omega=\sum t^a\otimes t_a\in
\mathfrak{sl}_3\otimes\mathfrak{sl}_3$
pairs each surviving direction with a BRST-exact direction
in the $f_{\min}$-reduction.  Explicitly:
$E_{23}\otimes E_{32}$ pairs the surviving root vector
$E_{23}$ (which gives $\Gp$) with $E_{32}$, which is in the
image of $\mathrm{ad}(f_{\min})$ and hence BRST-exact.
The projection $p\otimes p$ kills this cross-term.
The same argument applies to $\Gm$ (from $E_{31}$)
and $J$ (from the Cartan subalgebra, where ghost
cross-terms are killed by the ghost-number projection ---
see Remark~\ref{rem:bp-coproduct-resolved} for details).

Item (d): the stress tensor $T$ is a \emph{composite},
constructed as a Sugawara-type expression from the affine
currents plus ghost corrections.  The $\mathfrak{sl}_3$
Casimir has a surviving Cartan component
$H_1\otimes H_1$ (where $H_1$ is the Cartan element
dual to the surviving $U(1)$ direction).  Under the
projection $p\otimes p$, this maps to $J\otimes J$.
The spectral parameter $1/z$ arises from the propagator
in the HPL transfer formula: the Casimir component
contributes at first order in $1/z$.
\end{proof}

\begin{remark}[Principal vs.\ non-principal coproduct
dichotomy]
\label{rem:bp-coproduct-dichotomy}
\index{coproduct!principal vs non-principal dichotomy|textbf}
This reveals a clean structural dichotomy
(extending Remark~\ref{rem:bp-coproduct-resolved}
and Theorem~\ref{thm:gravitational-primitivity}):
\begin{itemize}
\item \textbf{Principal} ($\mathrm{Vir}$,
$\mathcal{W}_N$):
\emph{all} generators and composites are primitive.
No weight-$1$ generators survive the DS reduction.
The ghost-number obstruction kills all cross-terms.
\item \textbf{Non-principal} ($\mathcal{B}^{k'}$
and generalisations):
\emph{free generators} ($\Gp$, $\Gm$, $J$) are
primitive, but the \emph{composite} stress tensor $T$
is non-primitive.  The cross-term $J\otimes J/z$
arises from the surviving weight-$1$ Casimir component.
\end{itemize}
The non-primitivity of $T$ is a \emph{consequence} of
the Sugawara construction and does not obstruct the HPL
transfer: the coalgebra structure is determined by the
coproduct on \emph{cogenerators} (which are the generators
of $\overline{\mathcal{B}}$), and these are all primitive.
The composite non-primitivity is a derived phenomenon ---
the first matter-coupled gravity example in the ordered
bar framework where this distinction is visible.
\end{remark}

\medskip
\noindent\textbf{Part~VI: Summary and comparison.}

\begin{center}
\renewcommand{\arraystretch}{1.15}
\begin{tabular}{lccc}
\textbf{Feature} & \textbf{Virasoro}
& \textbf{SCA} & \textbf{BP} \\
\hline
Generators & $1$ ($T$) & $4$ ($T,\Gp,\Gm,J$)
& $4$ ($T,\Gp,\Gm,J$) \\
Depth classes & $\mathbf{M}$ only
& $\mathbf{G}$/$\mathbf{L}$/$\mathbf{C}$/$\mathbf{M}$
& $\mathbf{G}$/$\mathbf{L}$/$\mathbf{C}$/$\mathbf{M}$ \\
$m_3(T,T,T)$ & Virasoro $m_3$
& Virasoro $m_3$ at $c_{\mathrm{SCA}}$
& Virasoro $m_3$ at $c_{\mathrm{BP}}$ \\
Quartic $\to$ matter & ---
& confined & confined \\
$\Delta_{z,1}(T)$ & primitive
& non-primitive ($J\otimes J/z$)
& non-primitive ($J\otimes J/z$) \\
$\Delta_{z,1}(\Gp)$ & ---
& primitive & primitive \\
$U(1)$ decomposition & absent
& present & present \\
Koszul dual & $\mathrm{Vir}_{26-c}$
& $\mathcal{B}^{-k-6}$ (conjectural)
& SCA${}_{k''}$ (conjectural)
\end{tabular}
\end{center}
The SCA and BP share the same generator spectrum and depth
stratification, but differ in OPE coefficients: the
$\Gp$--$\Gm$ bracket \eqref{eq:bp-GpGm} contains the
Sugawara composite $T+\frac{3}{2K}{:}JJ{:}$, whereas the
SCA bracket (Computation~\ref{comp:ordered-bar-n2-sca})
contains $T+\frac{1}{2}\partial J$ without a normal-ordered
composite.  This difference is visible in the $m_2$ component
of the bar differential and propagates to all $m_n$ for
$n\ge 3$ via the HPL iteration.  The two algebras are
thus distinguished at the level of the ordered bar complex,
as expected for a Koszul-dual pair on opposite sides of
the complementarity identity~\eqref{eq:bp-complementarity}.
\end{computation}

\subsection{Modular extension: the ordered bar complex at
genus $g\ge 1$}
\label{subsec:ordered-bar-modular}

\begin{construction}[Modular ordered bar complex;
\ClaimStatusProvedHere]
\label{constr:modular-ordered-bar}
\index{bar complex!ordered!modular extension|textbf}
\index{modular R-matrix|textbf}
At genus~$g\ge 1$, the ordered configuration space on a curve
$\Sigma_g$ acquires new topology: the surface pure braid group
$\pi_1(\mathrm{Conf}_n^{\mathrm{ord}}(\Sigma_g))$ has generators
from the braid group (as at genus~$0$) and from the handle
loops of~$\Sigma_g$.  The monodromy of the ordered bar connection
around the handle loops produces genus corrections to the
$R$-matrix.

The \emph{modular ordered bar complex} is
\[
{\Barch}^{\mathrm{ord,mod}}(\cA)(g,n)
\;=\;
\bigoplus_{\Gamma\in\mathsf{Gr}^{\mathrm{st,ord}}_{g,n}}
\frac{\hbar^{g(\Gamma)}}{|\operatorname{Aut}(\Gamma)|}
\;(s^{-1}\bar\cA)^{\otimes|V(\Gamma)|}
\;\otimes\;
C_\bullet\!\bigl(
\mathrm{Conf}^{\mathrm{ord}}_{V(\Gamma)}(\Sigma_g)\bigr)
\;\otimes\;\det(E(\Gamma)),
\]
where $\mathsf{Gr}^{\mathrm{st,ord}}_{g,n}$ denotes stable
graphs whose vertices carry a \emph{planar} (total) ordering.

The \emph{modular $R$-matrix} is the genus expansion
\begin{equation}\label{eq:modular-r-matrix}
R^{\mathrm{mod}}(z;\hbar)
\;=\;
\sum_{g\ge 0}\hbar^{2g}\,r_g(z),
\end{equation}
with $r_0(z)=R(z)$ the genus-$0$ $R$-matrix and $r_g(z)$
($g\ge 1$) the genus-$g$ correction from the period structure
of~$\Sigma_g$.

At genus~$1$: the correction $r_1(z)$ arises from the monodromy
of the ordered bar connection around the $A$- and $B$-cycles
of the elliptic curve.  For
$\widehat{\mathfrak{sl}}_2$ at level~$k$:
\[
r_1(z)
\;=\;
\frac{\kappa}{(k+2)^2}\,
\frac{\wp(z;\tau)}{z}\,\Omega
\;+\;\text{(non-Casimir terms)},
\]
where $\wp$ is the Weierstrass function of the elliptic
curve with modulus~$\tau$, and $c=3k/(k+2)$ is the central charge.
The Weierstrass function appears because the genus-$1$
propagator $\varpi_{12}=\partial_{z_1}\partial_{z_2}\log
\theta_1(z_1-z_2;\tau)$ has an expansion
$\varpi_{12}=1/(z_1-z_2)^2-\wp(z_1-z_2;\tau)+\cdots$
whose non-rational part is~$\wp$.  The genus-$0$
$R$-matrix $R(z)=1+\Omega/z+\cdots$ acquires an
elliptic deformation $R_\tau(z)=1+\Omega/z
+\hbar^2\kappa\,\wp(z;\tau)\,\Omega/(k+2)^2+\cdots$.

For $\cH_k$ (Heisenberg): $r_1(z)=k\,\wp(z;\tau)/z$
(scalar).  The modular $R$-matrix is
$R^{\mathrm{mod}}(z;\hbar)=z^k+\hbar^2 k\,\wp(z;\tau)/z+\cdots$:
the Heisenberg braiding acquires elliptic corrections at genus~$1$.
\end{construction}

\subsection{The quantum lattice VOA: a genuinely $\Eone$-chiral algebra}%
\label{subsec:quantum-lattice-voa}%
\index{quantum lattice VOA|textbf}%
\index{E1-chiral algebra@$\Eone$-chiral algebra!quantum lattice VOA}%
\index{quantum vertex algebra!lattice type}

The preceding examples are all $E_\infty$-chiral (local vertex
algebras).  We now exhibit the first example of a
\emph{genuinely} $\Eone$-chiral algebra: the quantum
lattice VOA $V_q$ of Etingof--Kazhdan type, where the
$R$-matrix is independent input (not derivable from a
local OPE) and the factorisation structure is defined on
ordered configurations only.

\begin{computation}[Quantum lattice VOA $V_q$: genuinely
$\Eone$-chiral; \ClaimStatusConjectured]%
\label{comp:quantum-lattice-voa}%
\index{quantum lattice VOA!bar complex}%
\index{quantum lattice VOA!R-matrix}%
\index{quantum lattice VOA!Koszul dual}%
\index{quantum group!from quantum lattice VOA}%
The quantum lattice VOA $V_q$ (for $q = e^{i\pi\hbar}$
not a root of unity) is the $\Eone$-chiral deformation
of the $A_1$ lattice VOA $V_{\sqrt{2}\Z}$.  It is the
simplest example in which the ordered bar complex and
the $\Eone$-chiral framework are genuinely necessary:
the algebra is \emph{not} a vertex algebra in the BD
sense, and its factorisation structure does not descend
from unordered configurations.

\smallskip
\noindent
\emph{Generators and shifted OPE.}
$V_q$ has generators $J(z)$ (Cartan current, weight~$1$)
and $E^\pm(z)$ (vertex operators, weight~$1$), with the
\emph{quantum vertex algebra} OPE:
\begin{align}
J(z)\,J(w) &\;\sim\;
\frac{1}{(z-w)^2}\,,
\label{eq:qlattice-JJ}
\\[3pt]
J(z)\,E^\pm(w) &\;\sim\;
\frac{\pm 2\, E^\pm(w)}{z-w}\,,
\label{eq:qlattice-JE}
\\[3pt]
E^+(z)\,E^-(w) &\;\sim\;
\frac{1}{(z-w+\hbar)(z-w-\hbar)}\,
\bigl(q\,K^+(w) - q^{-1}\,K^-(w)\bigr),
\label{eq:qlattice-EE}
\end{align}
where $K^\pm(w) = q^{\pm J_0}$ are the quantum group
Cartan elements.  The critical feature is the
\emph{shifted poles} in~\eqref{eq:qlattice-EE}: the OPE
is singular at $z - w = \pm\hbar$, not at $z = w$.  This
shifting is the hallmark of a quantum vertex algebra:
the singularity locus depends on the deformation
parameter~$\hbar$.

\smallskip
\noindent
\emph{Class $\mathbf{L}$, depth~$3$.}
The shifted double pole in~\eqref{eq:qlattice-EE}
contributes to depth~$1$ (the quantum Killing form).
The shifted simple-pole residues contribute the quantum
Lie bracket (the $q$-commutator $[E^+, E^-]_q$) at
depth~$2$.  The quartic operation $m_4 = 0$ by the
quantum Jacobi identity (the quantum Serre relation
$(E^\pm)^2 E^\mp - (q + q^{-1})\,E^\pm E^\mp E^\pm
+ E^\mp (E^\pm)^2 = 0$).  Class~$\mathbf{L}$, depth~$3$.

\smallskip
\noindent
\emph{Operations: $m_2$ table.}
\begin{center}
\small
\renewcommand{\arraystretch}{1.3}
\begin{tabular}{lll}
\toprule
\textbf{Pair} & $m_2(a,b;\lambda)$ & \textbf{Source} \\
\midrule
$(J, J)$ & $\lambda$
  & double pole \\
$(J, E^\pm)$ & $\pm 2\, E^\pm$
  & simple pole \\
$(E^+, E^-)$
  & $\displaystyle\frac{qK^+ - q^{-1}K^-}{q - q^{-1}}$
  & shifted poles at $\pm\hbar$ \\
\bottomrule
\end{tabular}
\end{center}
The $m_2$ operation on $(E^+, E^-)$ differs from the
classical lattice VOA: the quantum group Cartan
elements $K^\pm$ replace the classical Cartan current~$J$,
and the denominator $q - q^{-1}$ is the quantum
normalisation.

\smallskip
\noindent
\emph{$R$-matrix: independent input (tier~(iii)).}
The $R$-matrix of $V_q$ is the \emph{trigonometric}
$R$-matrix of $U_q(\widehat{\mathfrak{sl}}_2)$:
\begin{equation}\label{eq:qlattice-R}
R(z)
\;=\;
\frac{1}{z - q^2}
\begin{pmatrix}
z\,q - q^{-1} & 0 & 0 & 0 \\
0 & z - 1 & q(z-1)(q - q^{-1})/z & 0 \\
0 & (q - q^{-1}) & z - 1 & 0 \\
0 & 0 & 0 & z\,q - q^{-1}
\end{pmatrix},
\end{equation}
where $z = e^{2\pi i u}$ is the multiplicative spectral
parameter.  This $R$-matrix is \emph{not} derivable from
the OPE by analytic continuation: the shifted poles at
$z = \pm\hbar$ produce the trigonometric dependence, and
the monodromy around $z = 0$ gives the quantum group
parameter~$q$.  This is tier~(iii) of the three-tier
picture: the $R$-matrix is independent input, defining
the nonlocal braiding.

\smallskip
\noindent
\emph{Koszul dual: $U_q(\widehat{\mathfrak{sl}}_2)$.}
The Koszul dual of $V_q$ is conjectured to be the
quantum affine algebra:
\begin{equation}\label{eq:qlattice-koszul}
V_q^!
\;\simeq\;
U_q(\widehat{\mathfrak{sl}}_2)\,.
\end{equation}
This is the $q$-deformation of the classical result
$V_{\sqrt{2}\Z}^! = Y_\hbar(\mathfrak{sl}_2)$
(Computation~\ref{comp:lattice-voa-ordered-bar}).
The rational $R$-matrix $R(z) = 1 + \Omega/z$
(Yang $R$-matrix) is deformed to the trigonometric
$R$-matrix~\eqref{eq:qlattice-R}, and the Yangian
$Y_\hbar(\mathfrak{sl}_2)$ is deformed to the quantum
affine algebra $U_q(\widehat{\mathfrak{sl}}_2)$.
The identification~\eqref{eq:qlattice-koszul} is
supported by the RTT presentation: the RTT relation
$R(z/w)\,T_1(z)\,T_2(w) = T_2(w)\,T_1(z)\,R(z/w)$
with the trigonometric $R$-matrix~\eqref{eq:qlattice-R}
generates $U_q(\widehat{\mathfrak{sl}}_2)$.

\smallskip
\noindent
\emph{Poincar\'{e} series.}
$P(t) = 1 + 3t$ (three cogenerators $s^{-1}J$,
$s^{-1}E^+$, $s^{-1}E^-$, matching the
$\dim(\mathfrak{sl}_2) = 3$ generators).

\smallskip
\noindent
\emph{Euler characteristic and modular data.}
The modular characteristic $\kappa(V_q) = 3q^2/(q^2 + q^{-2} + 1)$
interpolates between the classical value $\kappa = 3/4$
(at $q = 1$, the $A_1$ lattice VOA at $k = 1$) and the
critical value $\kappa = 0$ (at $q = i$, a root of unity).

\smallskip
\noindent
\emph{Genuinely $\Eone$: the three diagnostic tests.}
Three properties distinguish $V_q$ from
$E_\infty$-chiral algebras:
\begin{enumerate}[label=\textup{(\roman*)},nosep]
\item \emph{Shifted singularities}: the OPE~\eqref{eq:qlattice-EE}
  has poles at $z - w = \pm\hbar$, not at $z = w$.  No local
  (unordered) factorisation structure is compatible with
  $\hbar$-shifted singularity loci.
\item \emph{Non-derived $R$-matrix}: the trigonometric
  $R$-matrix~\eqref{eq:qlattice-R} is not the monodromy
  of any flat connection on $\Conf_2(\C)$ (it has a
  $q$-dependent spectral parameter, not a rational one).
\item \emph{Ordered-only factorisation}: the $n$-point
  correlation functions are defined on
  $\Conf_n^{\mathrm{ord}}(\C)$ and the Yang--Baxter
  equation replaces the $\Sigma_n$-equivariance of
  the $E_\infty$ setting.
\end{enumerate}
This is the mathematical content of ``nonlocal'' in the
$\Eone$-chiral framework: the braiding data cannot be
recovered from local OPE data, and the ordered configuration
space is the primitive geometric input.

\smallskip
\noindent
\emph{Genus-$1$.}
Propagator-entangled ($c_0 \ne 0$).  The genus-$1$
deformation produces the \emph{elliptic quantum group}
$U_{q,p}(\widehat{\mathfrak{sl}}_2)$ of
Felder--Varchenko: the trigonometric $R$-matrix acquires
an elliptic deformation
$R(z;\tau) = R_{\mathrm{trig}}(z) + O(q_\tau)$,
where $q_\tau = e^{2\pi i \tau}$ is the elliptic nome.
The passage rational $\to$ trigonometric $\to$ elliptic
mirrors the genus tower: genus~$0$ gives the Yangian,
the $q$-deformation gives the quantum affine algebra,
and genus~$1$ gives the elliptic quantum group.
\ClaimStatusConjectured
\end{computation}

\begin{remark}[The $\Eone$ landscape: from vertex algebras to quantum vertex algebras]%
\label{rem:e1-landscape-quantum}%
\index{E1-chiral algebra@$\Eone$-chiral algebra!landscape}%
\index{quantum vertex algebra!landscape}%
Computation~\ref{comp:quantum-lattice-voa} provides the
first fully worked example in the genuinely $\Eone$ column
of the ordered bar landscape.  The classical ($q = 1$) limit
recovers the $A_1$ lattice VOA
(Computation~\ref{comp:lattice-voa-ordered-bar}), and the
three diagnostic tests isolate the features that make the
quantum deformation genuinely nonlocal.  The conjectured
Koszul dual identification $V_q^! \simeq
U_q(\widehat{\mathfrak{sl}}_2)$ (equation~\eqref{eq:qlattice-koszul})
extends the rational Koszul duality
$V_\Lambda^! = Y_\hbar(\fg)$ to the trigonometric regime.
Proving~\eqref{eq:qlattice-koszul} would establish the
full Yang--Baxter hierarchy
(rational $\to$ trigonometric $\to$ elliptic) as a
manifestation of the ordered bar complex at successive
genus levels.
\end{remark}
